\renewcommand{\chaptername}{Κεφάλαιο}
\chapter{Εκτεταμένη Περίληψη στα Ελληνικά}
\label{chapter:greek}
\renewcommand{\tablename}{Πίνακας}
\renewcommand{\figurename}{Σχήμα}
\clearpage

\section{Εισαγωγή}

Τα Μεγάλα Γλωσσικά Μοντέλα (ΜΓΜ - Large Language Models) ορίζονται ως παραμετρικά μοντέλα πιθανοτικής γλώσσας μεγάλης κλίμακας, βασισμένα στην αρχιτεκτονική Transformer. Τα μοντέλα αυτά εκπαιδεύονται μέσω αντικειμενικών συναρτήσεων αυτο-επίβλεψης για την εκτίμηση της κατανομής πιθανότητας εκτεταμένων ακολουθιών φυσικής γλώσσας. Μέσα από αυτή τη διαδικασία εκμάθησης σε ετερογενή σώματα κειμένου (corpora), τα ΜΓΜ αναπτύσσουν αναδυόμενες γνωστικές ικανότητες (emergent abilities), οι οποίες τους επιτρέπουν να παράγουν συνεκτικό λόγο, να επιλύουν σύνθετα προβλήματα μέσω πολυβηματικού συλλογισμού (chain-of-thought reasoning) και να εξάγουν σημασιολογικά πλούσιες αναπαραστάσεις.

Παρά τη δομική τους επάρκεια, τα ΜΓΜ περιορίζονται εγγενώς από τη στατικότητα της παραμετρικής τους γνώσης, η οποία κρυσταλλώνεται κατά τη φάση της προεκπαίδευσης. Η επικαιροποίηση αυτής της γνώσης ή η ενσωμάτωση εξειδικευμένων πληροφοριών απαιτεί υπολογιστικά δαπανηρές διαδικασίες επανεκπαίδευσης ή λεπτομερούς ρύθμισης (fine-tuning). Επιπλέον, οι παραγόμενες αποκρίσεις στερούνται ενός εγγενούς μηχανισμού ανιχνευσιμότητας (traceability) και επαλήθευσης από εξωτερικές πηγές. Το γεγονός αυτό οδηγεί συχνά στην εκδήλωση ανυπόστατων πραγματολογικών σφαλμάτων, γνωστών ως ψευδαισθήσεις (hallucinations), περιορίζοντας την αξιοπιστία των μοντέλων σε πεδία υψηλής κρισιμότητας.

Ως απάντηση στις αδυναμίες αυτές, η Ανάκτηση-Επαυξημένη Παραγωγή (Retrieval-Augmented Generation - RAG) εισάγει ένα υβριδικό παραδειγματικό πλαίσιο, το οποίο συζευγνύει την παραμετρική μνήμη του γλωσσικού μοντέλου με μη παραμετρικά, εξωτερικά αποθετήρια γνώσης. Ένα τυπικό σύστημα RAG αποτελείται από δύο βασικές συνιστώσες: έναν ανακτητή (retriever), ο οποίος χαρτογραφεί ένα ερώτημα χρήστη στα πιο σχετικά αποσπάσματα κειμένου ενός σώματος εγγράφων, και έναν παραγωγό (generator), ο οποίος συνθέτει την τελική απόκριση λαμβάνοντας υπόψη τα ανακτηθέντα συμφραζόμενα. Η αρχιτεκτονική αυτή ενισχύει την πραγματολογική ακρίβεια, επιτρέπει τη δυναμική προσαρμογή σε μεταβαλλόμενα δεδομένα και προσφέρει τη δυνατότητα παράθεσης πηγών.

Ωστόσο, η τρέχουσα βιβλιογραφία εστιάζει κατά κύριο λόγο σε σενάρια μονοστροφικής αλληλεπίδρασης (single-turn RAG), όπου κάθε ερώτημα αντιμετωπίζεται ως απομονωμένο και ανεξάρτητο. Η υπόθεση αυτή αποκλίνει από τα πραγματικά περιβάλλοντα χρήσης, τα οποία είναι εγγενώς διαλογικά και πολύστροφα (multi-turn). Στα συστήματα Πολύστροφου RAG (Multi-Turn RAG), οι διαδοχικές ερωτήσεις χαρακτηρίζονται από έντονες διαλογικές εξαρτήσεις, όπως η χρήση αντωνυμιών χωρίς ρητό υποκείμενο (coreference resolution) και η αποσπασματική διατύπωση ερωτημάτων. Η παράλειψη αυτών των εξαρτήσεων προκαλεί μετατόπιση των συμφραζομένων (context drift), έντονη αστάθεια στην ανάκτηση και σωρευτική διάδοση σφαλμάτων (error cascade) κατά την εξέλιξη του διαλόγου.

Στην παρούσα διπλωματική εργασία, το πρόβλημα του Πολύστροφου RAG προσεγγίζεται υπό το πρίσμα της \textit{σταθερότητας ανάκτησης συμφραζομένων} (context-aware retrieval stability). Υποστηρίζουμε ότι οι μικρές αποκλίσεις κατά την αναδιατύπωση των ερωτημάτων, σε συνδυασμό με το χάσμα λεξιλογίου (vocabulary gap) μεταξύ των στροφών, προκαλούν απότομη υποβάθμιση της ποιότητας των ανακτηθέντων εγγράφων μετά την πρώτη στροφή. Για την επίλυση αυτής της πρόκλησης, σχεδιάστηκε και υλοποιήθηκε μια ενοποιημένη αρχιτεκτονική που διαχωρίζει λειτουργικά τις εξής κρίσιμες διεργασίες:
\begin{enumerate}
    \item[$(\alpha)$] Την αναπαράσταση και εμπλουτισμό του διαλογικού ιστορικού.
    \item[$(\beta)$] Την ανακάλυψη και δυναμική σύνθεση αποδεικτικών στοιχείων μέσω πολλαπλών στρατηγικών.
    \item[$(\gamma)$] Τη λήψη απόφασης για την τελική παραγωγή μέσω ελεγχόμενων πυλών απαντησιμότητας (answerability gates).
\end{enumerate}

Η σταθεροποίηση του μηχανισμού ανάκτησης επιτυγχάνεται μέσω μιας προσέγγισης \textit{ποικιλότητας ερωτημάτων έναντι ποικιλότητας ανακτητών} (query-diversity-over-retriever-diversity). Αντί της χρήσης ετερογενών ανακτητών, το σύστημα εκδίδει πέντε συμπληρωματικές αναδιατυπώσεις ερωτήματος βασισμένες σε ΜΓΜ προς έναν ενιαίο, ευθυγραμμισμένο με το σώμα κειμένου (corpus-aligned) αραιό ανακτητή (sparse retriever). Η σύνθεση των αποτελεσμάτων πραγματοποιείται μέσω ενός νέου αλγορίθμου Ιεραρχικής Συγχώνευσης Αντίστροφης Κατάταξης (Nested Reciprocal Rank Fusion - Nested RRF), ο οποίος συγκεντρώνει πρώτα τις πιο διερευνητικές στρατηγικές σε μια ενδιάμεση κατάταξη και στη συνέχεια τη συνδυάζει με τις πιο συντηρητικές. Παράλληλα, ο μηχανισμός παραγωγής βασίζεται σε μια στρατηγική \textit{καθυστερημένης δέσμευσης} (deferred commitment), η οποία αποσυνθέτει τη διαδικασία σε εξαγωγή αποσπασμάτων (evidence span extraction), διαμόρφωση εναλλακτικών απαντήσεων και τελική επιλογή κατόπιν συμφωνίας ενός συστήματος πολλαπλών κριτών (multi-judge selection).

Η προτεινόμενη αρχιτεκτονική αξιολογήθηκε στο πλαίσιο του διεθνούς διαγωνισμού SemEval-2026 Task 8 (MTRAGEval), ο οποίος εξετάζει τρεις διασυνδεδεμένες εργασίες: την ανάκτηση αποσπασμάτων (Task A), την παραγωγή απάντησης με βάση δεδομένες πηγές (Task B) και την end-to-end υλοποίηση συστήματος RAG (Task C). Το αναπτυχθέν σύστημα, με την ονομασία \textbf{AILS-NTUA}, πέτυχε υψηλές επιδόσεις στις δύο πρώτες εργασίες και κατέλαβε την \textbf{1η θέση στο Task A} (nDCG@5: 0,5776, βελτίωση +20,5\% έναντι της ισχυρότερης γραμμής βάσης, μεταξύ 38 συμμετοχών) και τη \textbf{2η θέση στο Task B} (Αρμονικός Μέσος - HM: 0,7698, μεταξύ 26 συστημάτων). Στο πιο απαιτητικό end-to-end σενάριο (\textbf{Task C}), το σύστημα κατατάχθηκε 11ο μεταξύ 29 συμμετοχών (HM: 0,5409), μόλις οριακά πάνω από την ισχυρότερη γραμμή βάσης των διοργανωτών (0,5366). Όπως αναλύεται στην Ενότητα 1.5, καθοριστικός παράγοντας της υστέρησης είναι η βαθμονόμηση της πύλης απαντησιμότητας: η πύλη βαθμονομήθηκε στο development set, ενώ το ποσοστό των μη απαντήσιμων ερωτημάτων στο test set είναι περίπου τριπλάσιο. Συμβάλλουν επίσης τα σφάλματα ανάκτησης στις κορυφαίες θέσεις, αφού στο Task C ο generator λαμβάνει μόνο τα τρία πρώτα αποσπάσματα που ανακτά το ίδιο το σύστημα.

Τα κύρια ερευνητικά ευρήματα της εργασίας συνοψίζονται στα ακόλουθα:
\begin{enumerate}
    \item Η εισαγωγή στοχευμένης ποικιλότητας ερωτημάτων (query diversity) σε έναν κατάλληλα ευθυγραμμισμένο ανακτητή υπερτερεί, στο development set και με βάση τις επίσημες κρίσεις συνάφειας του benchmark, της σύνθεσης πολλών διαφορετικών μοντέλων ανάκτησης, καθώς περιορίζει το vocabulary gap χωρίς να εισάγει τον θόρυβο της σύντηξης ετερογενών κατατάξεων.
    \item Οι μηχανισμοί επαναβαθμολόγησης (cross-encoders) λειτουργούν βέλτιστα ως συμπληρωματικά φίλτρα διύλισης και δεν μπορούν να υποκαταστήσουν την ανάγκη για ένα σταθερό αρχικό στάδιο ανάκτησης.
    \item Η χρήση ΜΓΜ για την περαιτέρω επαναβαθμολόγηση των αποτελεσμάτων παρουσιάζει φαινόμενα κορεσμού (saturation) όταν η ποιότητα της αρχικής ανάκτησης είναι ήδη υψηλή, αυξάνοντας παράλληλα το υπολογιστικό κόστος χωρίς αξιόλογο όφελος.
    \item Καθοριστικός παράγοντας της end-to-end απόδοσης (Task C) υπό την κατανομή του test set είναι η βαθμονόμηση της πύλης απαντησιμότητας (answerability gate), ενώ τα σφάλματα ανάκτησης στις κορυφαίες θέσεις παραμένουν σημαντικά. Η πύλη επιβαρύνεται από την τάση των μοντέλων να αποδέχονται αδύναμες αποδείξεις (acceptance bias) και γίνεται κρίσιμη όταν αυξάνεται το ποσοστό των μη απαντήσιμων ερωτημάτων.
\end{enumerate}

\paragraph{Σχέση με δημοσιευμένη εργασία και προσωπική συνεισφορά.}
Το σύστημα αναπτύχθηκε στην ομάδα AILS-NTUA και δημοσιεύτηκε ως άρθρο περιγραφής συστήματος στο SemEval-2026~\citep{athanasiou2026ailsntua}· η παρούσα διπλωματική βασίζεται σε αυτό και επεκτείνει την ανάλυσή του. Ο συγγραφέας σχεδίασε και υλοποίησε τις συνιστώσες ανάκτησης, παραγωγής και απαντησιμότητας, εκτέλεσε όλα τα πειράματα και τις αναλύσεις και έγραψε τη διπλωματική· η Μ.~Λυμπεραίου και ο Γ.~Φιλανδριανός προσέφεραν ερευνητική καθοδήγηση και ανατροφοδότηση, ενώ οι Α.~Βουλόδημος και Γ.~Στάμου επόπτευσαν την εργασία.

%==============================================================
\section{Θεωρητικό Υπόβαθρο}
%==============================================================

Η παρούσα ενότητα παρουσιάζει τις θεμελιώδεις θεωρητικές έννοιες και το μαθηματικό υπόβαθρο που υποστηρίζουν τη σχεδίαση συστημάτων Ανάκτησης-Επαυξημένης Παραγωγής (Retrieval-Augmented Generation - RAG) σε πολύστροφα (multi-turn) διαλογικά περιβάλλοντα. Συγκεκριμένα, αναλύεται η στοχαστική φύση των Μεγάλων Γλωσσικών Μοντέλων, το υβριδικό παραδειγματικό πλαίσιο της αρχιτεκτονικής RAG, οι σύγχρονες μέθοδοι νευρωνικής ανάκτησης πληροφορίας, οι τεχνικές αναδιατύπωσης ερωτημάτων μέσω συλλογιστικής, καθώς και οι προκλήσεις της συνομιλιακής αναζήτησης και οι μεθοδολογίες πολυδιάστατης αξιολόγησής τους.

%--------------------------------------------------------------
\subsection{Μεγάλα Γλωσσικά Μοντέλα}
%--------------------------------------------------------------

Τα Μεγάλα Γλωσσικά Μοντέλα (ΜΓΜ - Large Language Models) αποτελούν αυτο-παλινδρομικά (autoregressive) στοχαστικά μοντέλα γλώσσας, βασισμένα στην αρχιτεκτονική Transformer και ειδικότερα στον μηχανισμό της αυτο-προσοχής (self-attention). Τα μοντέλα αυτά εκπαιδεύονται μέσω μεγιστοποίησης της λογαριθμικής πιθανότητας του επόμενου συμβόλου (next-token prediction) πάνω σε εκτεταμένα σώματα κειμένου, μια διαδικασία που τυπικά μοντελοποιείται ως:

\begin{equation}
P(X) = \prod_{t=1}^{T} P(x_t \mid x_1, x_2, \dots, x_{t-1}) = \prod_{t=1}^{T} P(x_t \mid x_{<t})
\end{equation}

Η εκθετική κλιμάκωση (scaling laws) του αριθμού των παραμέτρων, του όγκου των δεδομένων εκπαίδευσης και της υπολογιστικής ισχύος οδηγεί στην εμφάνιση αναδυόμενων ικανοτήτων (emergent abilities). Οι σημαντικότερες από αυτές περιλαμβάνουν:
\begin{itemize}
    \item \textbf{Μάθηση εντός πλαισίου (In-Context Learning - ICL):} Η ικανότητα εκτέλεσης νέων εργασιών κατά τη φάση της εξαγωγής συμπερασμάτων (inference) αποκλειστικά μέσω της παροχής παραδειγμάτων στην προτροπή (prompt), χωρίς τροποποίηση των βαρών του δικτύου.
    \item \textbf{Συμμόρφωση με οδηγίες (Instruction Following):} Η ευθυγράμμιση του μοντέλου με τις προθέσεις του χρήστη μέσω εκπαίδευσης εποπτευόμενης λεπτομερούς ρύθμισης (Supervised Fine-Tuning - SFT) και ενισχυτικής μάθησης από ανθρώπινη ανάδραση (RLHF).
    \item \textbf{Πολυβηματικός Συλλογισμός (Chain-of-Thought Reasoning):} Η ικανότητα αποδόμησης σύνθετων προβλημάτων σε διαδοχικά, λογικά ενδιάμεσα βήματα συλλογιστικής πριν από την παραγωγή της τελικής απόκρισης. Στο Κεφάλαιο 5 αξιοποιούμε άμεσα αυτή την αρχή στη σχεδίαση της στρατηγικής Chain-of-Thought αναδιατύπωσης ερωτήματος (Ενότητα 5.2.1) και στην πολυσταδιακή διαδικασία παραγωγής απάντησης (Ενότητα 5.3).
\end{itemize}

Παρά την υψηλή τους αναπαραστατική ικανότητα, τα ΜΓΜ παρουσιάζουν δομικούς περιορισμούς:
\begin{itemize}
    \item \textbf{Στατικότητα γνώσης (Knowledge Cutoff):} Η ενσωματωμένη γνώση παραμένει χρονικά και πραγματολογικά περιορισμένη στο σημείο ολοκλήρωσης της προεκπαίδευσης.
    \item \textbf{Έλλειψη εγγενούς ανιχνευσιμότητας (Traceability):} Τα μοντέλα λειτουργούν ως "μαύρα κουτιά" (black boxes), αδυνατώντας να παρέχουν ρητές παραπομπές ή αποδείξεις για την ορθότητα των ισχυρισμών τους.
    \item \textbf{Πραγματολογικές Ψευδαισθήσεις (Hallucinations):} Η τάση παραγωγής κειμένου που είναι συντακτικά και υφολογικά άρτιο, αλλά πραγματολογικά εσφαλμένο ή ανυπόστατο, λόγω της φύσης τους ως προσεγγιστές κατανομών πιθανότητας.
\end{itemize}

Οι περιορισμοί αυτοί επιβάλλουν τη μετάβαση από αμιγώς παραμετρικά συστήματα σε υβριδικές αρχιτεκτονικές που ενσωματώνουν εξωτερικές, δυναμικές πηγές γνώσης.

%--------------------------------------------------------------
\subsection{Μέθοδοι Ανάκτησης Πληροφορίας}
%--------------------------------------------------------------

Οι σύγχρονες αρχιτεκτονικές ανάκτησης πληροφορίας κατηγοριοποιούνται με βάση τον τρόπο αναπαράστασης του κειμένου:

\begin{itemize}
    \item \textbf{Λεξιλογικοί Ανακτητές (Sparse Retrievers):} Βασίζονται στην ακριβή αντιστοίχιση όρων (exact keyword matching) χρησιμοποιώντας στατιστικές συχνότητας εμφάνισης. Ο αλγόριθμος BM25 αποτελεί το τυπικό βιομηχανικό πρότυπο, βαθμολογώντας τη συνάφεια ενός εγγράφου $d$ με βάση τον όρο $t$ ως συνάρτηση της συχνότητας του όρου εντός του εγγράφου (term frequency) και της αντίστροφης συχνότητας εμφάνισης στο σώμα κειμένου (inverse document frequency). Παρά την υψηλή ακρίβεια σε τεχνικούς όρους και σειριακούς αριθμούς, το BM25 περιορίζεται θεμελιωδώς από το πρόβλημα της ασυμφωνίας λεξιλογίου (vocabulary mismatch): αν το ερώτημα χρησιμοποιεί παραφράσεις ή συνώνυμα που δεν συμπίπτουν λεξιλογικά με το κείμενο-στόχο, η τομή $q \cap d$ μηδενίζεται.
    \item \textbf{Πυκνοί Ανακτητές (Dense Retrievers):} Χρησιμοποιούν βαθιά νευρωνικά δίκτυα (π.χ. αρχιτεκτονικές Dual-Encoder) εκπαιδευμένα μέσω αντιπαραθετικής μάθησης (contrastive learning) για τη χαρτογράφηση ερωτημάτων και εγγράφων σε έναν κοινό, χαμηλής διάστασης διανυσματικό χώρο $\mathbb{R}^d$. Η βαθμολογία συνάφειας υπολογίζεται μέσω του εσωτερικού γινομένου ή της συνημίτονου ομοιότητας των αντίστοιχων embeddings: $S(q, d) = \langle \mathbf{E}_Q(q), \mathbf{E}_D(d) \rangle$. Παρότι αποτελεσματικοί σε αφηρημένη σημασιολογική ταύτιση και διαγλωσσικές παραφράσεις, εισάγουν συχνά σημασιολογικό θόρυβο σε πολύστροφα περιβάλλοντα, υπεργενικεύοντας τεχνικά tokens.
    \item \textbf{Μαθητοί Αραιοί Ανακτητές (Learned Sparse Retrievers):} Αποτελούν μια σύγχρονη υβριδική προσέγγιση (π.χ. SPLADE, ELSER) που ενοποιεί τα πλεονεκτήματα λεξιλογικής ακρίβειας και σημασιολογικής γενίκευσης. Χρησιμοποιούν ένα γλωσσικό μοντέλο για να προβλέψουν και να επεκτείνουν το λεξιλόγιο ενός εγγράφου (term expansion) στον αρχικό χώρο των high-dimensional tokens, επιλύοντας το πρόβλημα της συνωνυμίας χωρίς να χάνουν την ακρίβεια της αραιής αναπαράστασης. Η βαθμολογία ορίζεται ως:
    \begin{equation}
    S(q, d) = \sum_{t \in q \cap d} w_q(t) \cdot w_d(t)
    \end{equation}
    όπου τα βάρη $w(t)$ είναι μαθημένες μη αρνητικές τιμές που αντανακλούν τη σημασιολογική βαρύτητα του όρου εντός του συγκεκριμένου corpus.
\end{itemize}

Ο Πίνακας~\ref{tab:retrieval-comparison} συνοψίζει τα τρία αυτά παραδείγματα ως προς τα κρίσιμα χαρακτηριστικά τους για το πολύστροφο σενάριο. Η παρούσα εργασία βασίζεται σε έναν corpus-aligned learned sparse ανακτητή (\textbf{ELSER v1}), ο οποίος, όπως αναλύεται διεξοδικά στην Ενότητα 1.5, υπερτερεί εμπειρικά τόσο του BM25 όσο και των ισχυρότερων εμπορικών dense ανακτητών (π.χ. Cohere Embed v4) στο benchmark MTRAG, με το πλεονέκτημα να αποδίδεται εν μέρει στη μείωση του vocabulary gap χωρίς τον σημασιολογικό θόρυβο των πυκνών embeddings (με την επιφύλαξη της μεροληψίας των κρίσεων συνάφειας, Ενότητα 1.5).

\begin{table}[h]
\centering
\small
\begin{tabular}{lccc}
\toprule
\textbf{Χαρακτηριστικό} & \textbf{Sparse (BM25)} & \textbf{Dense} & \textbf{Learned Sparse (ELSER v1)} \\
\midrule
Αναπαράσταση & Λεξιλογική, αραιή & Νευρωνική, πυκνή ($\mathbb{R}^d$) & Λεξιλογική, αραιή με term expansion \\
Χειρισμός συνωνύμων & Όχι & Ναι & Ναι (μέσω επέκτασης όρων) \\
Ακρίβεια σε IDs/κωδικούς & Υψηλή & Χαμηλή & Υψηλή \\
Ερμηνευσιμότητα & Πλήρης & Περιορισμένη & Υψηλή \\
Ευθυγράμμιση με corpus & — & Γενικού σκοπού & Corpus-aligned \\
Επίδοση στο MTRAG (R@5) & 0.153 & 0.397 (καλύτερο dense) & \textbf{0.483} \\
\bottomrule
\end{tabular}
\caption{Συγκριτική επισκόπηση παραδειγμάτων ανάκτησης. Οι τιμές Recall@5 στη στήλη MTRAG προέρχονται από την αξιολόγηση χωρίς αναδιατύπωση ερωτήματος (777 ερωτήματα development set) και αιτιολογούν την επιλογή του ELSER v1 ως βάσης ανάκτησης του προτεινόμενου συστήματος.}
\label{tab:retrieval-comparison}
\end{table}

%--------------------------------------------------------------
\subsection{Αναδιατύπωση Ερωτημάτων και Σύντηξη}
%--------------------------------------------------------------

Σε δυναμικά διαλογικά περιβάλλοντα, τα ερωτήματα των χρηστών είναι συχνά ελλιπή ή εξαρτώνται άμεσα από το ιστορικό. Η Αναδιατύπωση Ερωτημάτων (Query Reformulation) μέσω ΜΓΜ μετασχηματίζει ένα εξαρτημένο ερώτημα σε μια αυτοτελή (context-independent) μορφή. Οι κυριότερες στρατηγικές περιλαμβάνουν:
\begin{itemize}
    \item \textbf{HyDE (Hypothetical Document Embeddings):} Παραγωγή μιας υποθετικής απάντησης (zero-shot) από το ΜΓΜ, η οποία στη συνέχεια χρησιμοποιείται ως διευρυμένο ερώτημα για την ανάκτηση, στοχεύοντας στην ευθυγράμμιση στον χώρο των απαντήσεων.
    \item \textbf{Chain-of-Thought (CoT) Reformulation:} Παραγωγή ενδιάμεσων βημάτων ανάλυσης της πληροφοριακής ανάγκης πριν την τελική εξαγωγή των όρων αναζήτησης.
    \item \textbf{Keyword Extraction \& Anchor Targeting:} Εστίαση στην εξαγωγή ονοματοποιημένων οντοτήτων και λέξεων-κλειδιών που λειτουργούν ως άγκυρες (anchors) στο συγκεκριμένο σώμα κειμένου.
\end{itemize}

Όταν εκδίδονται πολλαπλές, συμπληρωματικές αναδιατυπώσεις ερωτήματος, απαιτείται ένας μηχανισμός σύντηξης για τη συγχώνευση των επιμέρους κατατάξεων. Η Σύντηξη Αντίστροφης Κατάταξης (Reciprocal Rank Fusion - RRF) αποτελεί μια ανθεκτική μέθοδο (unsupervised rank aggregation) που ορίζεται ως:

\begin{equation}
\text{RRF\_Score}(d \in \mathcal{D}) = \sum_{m \in \mathcal{M}} \frac{1}{k + r_m(d)}
\end{equation}

όπου $\mathcal{M}$ είναι το σύνολο των παραχθέντων λιστών κατάταξης, $r_m(d)$ είναι η θέση (rank) του εγγράφου $d$ στη λίστα $m$, και $k$ είναι μια σταθερά εξομάλυνσης (smoothing constant, τυπικά $k=60$). Η RRF εξασφαλίζει ότι έγγραφα που κατατάσσονται ψηλά σε πολλαπλές στρατηγικές ευνοούνται, χωρίς να επηρεάζονται από τις αποκλίσεις στις απόλυτες τιμές των σκορ των διαφορετικών μοντέλων. Ωστόσο, η επίπεδη (flat) εκδοχή της RRF αντιμετωπίζει όλες τις στρατηγικές αναδιατύπωσης ως ισοδύναμες, κάτι που μπορεί να είναι προβληματικό όταν ορισμένες (π.χ. HyDE, CoT) αναμένεται να είναι λιγότερο σταθερές ανά στροφή από άλλες (π.χ. Minimal, Corpus-Specific). Μια ιεραρχική παραλλαγή της RRF που αντιμετωπίζει αυτή την ανομοιογένεια παρουσιάζεται στην Ενότητα 1.4.2.

%--------------------------------------------------------------
\subsection{Πολύστροφη Συνομιλιακή RAG}
%--------------------------------------------------------------

Σε αντίθεση με το μονοστροφικό RAG, η Πολύστροφη Συνομιλιακή RAG (Multi-Turn Conversational RAG) εισάγει τη διάσταση του χρόνου και της διαλογικής εξάρτησης. Το κύριο πρόβλημα εντοπίζεται στη διατήρηση της \textit{σταθερότητας ανάκτησης συμφραζομένων} (context-aware retrieval stability). Καθώς ο διάλογος εξελίσσεται, τα ερωτήματα εμφανίζουν:
\begin{itemize}
    \item \textbf{Επίλυση Αναφορών (Coreference Resolution):} Χρήση αντωνυμιών ή δεικτικών εκφράσεων που αναφέρονται σε οντότητες προηγούμενων στροφών.
    \item \textbf{Ελλειπτικότητα (Ellipsis):} Παράλειψη ολόκληρων συντακτικών δομών ή υποκειμένων που θεωρούνται γνωστά από τα συμφραζόμενα.
    \item \textbf{Μετατόπιση Θεματικού Πλαισίου (Topic Switching):} Απότομες αλλαγές στη ροή της συζήτησης.
\end{itemize}

Η μη επιτυχής διαχείριση αυτών των φαινομένων οδηγεί σε \textit{σωρευτική διάδοση σφαλμάτων} (error cascading): ένα αποτυχημένο βήμα ανάκτησης στη στροφή $t$ εισάγει θόρυβο στη γεννήτρια, η οποία παράγει μια εσφαλμένη απάντηση, η οποία με τη σειρά της μολύνει το ιστορικό του διαλόγου, προκαλώντας ακόμα μεγαλύτερη απόκλιση (context drift) στη στροφή $t+1$. Ο μηχανισμός αυτός προϋποθέτει ότι το ιστορικό περιέχει τις προηγούμενες απαντήσεις του ίδιου του συστήματος, όπως σε έναν πραγματικό βοηθό· το MTRAGEval, αντίθετα, παρέχει σε κάθε στροφή το ιστορικό αναφοράς, οπότε η αξιολόγησή μας δεν τον ενεργοποιεί. Στο MTRAG benchmark παρατηρούμε πάντως απότομη πτώση της απόδοσης ανάκτησης μετά την πρώτη στροφή (Ενότητα 1.5.2).

%--------------------------------------------------------------
\subsection{Αξιολόγηση Συστημάτων RAG}
%--------------------------------------------------------------

Η αξιολόγηση σε πολύστροφα περιβάλλοντα RAG επιβάλλει τη χρήση ενός πολυδιάστατου πλαισίου μετρικών που εξετάζει ανεξάρτητα την ανάκτηση και την παραγωγή. Οι βασικοί άξονες περιλαμβάνουν:
\begin{itemize}
    \item \textbf{Μετρικές Ανάκτησης εγγράφων:} Το Recall@k (ικανότητα εντοπισμού όλων των σχετικών εγγράφων) και το Normalized Discounted Cumulative Gain (nDCG@k), το οποίο λαμβάνει υπόψη τη σχετική θέση και τη διαβαθμισμένη συνάφεια των ανακτηθέντων αποσπασμάτων.
    \item \textbf{Μετρικές Πιστότητας και Παραγωγής:} Το αναφορικά αλγοριθμικό σκορ $RB_{alg}$ (σύνθεση BERT-Recall, BERT-K-Precision και ROUGE-L έναντι της ανθρώπινα επιμελημένης αναφοράς) και το χωρίς αναφορά μετρικό πιστότητας $RL_F$ (πλαίσιο RAGAS), που αποτιμά αν οι ισχυρισμοί της απάντησης θεμελιώνονται στα ανακτηθέντα αποσπάσματα.
    \item \textbf{Αξιολόγηση μέσω ΜΓΜ (LLM-as-a-Judge):} Το αναφορικό σκορ $RB_{llm}$, ένας διάμεσος-συγκεντρωμένος αξιολογητής πολλαπλών ΜΓΜ που εκτιμά Πιστότητα (Faithfulness), Καταλληλότητα (Appropriateness) και Πληρότητα (Completeness) της απάντησης.
\end{itemize}

Στο πλαίσιο του MTRAGEval (SemEval-2026 Task 8), η end-to-end απόδοση του συστήματος συντίθεται μέσω του Αρμονικού Μέσου (Harmonic Mean - HM) των επιμέρους αυτών διαστάσεων:

\begin{equation}
\text{HM} = \frac{3}{\frac{1}{\text{RB}_{alg}} + \frac{1}{\text{RL}_F} + \frac{1}{\text{RB}_{llm}}}
\end{equation}

Η χρήση του αρμονικού μέσου τιμωρεί αυστηρά συστήματα που παρουσιάζουν ασυμμετρία στην απόδοσή τους (π.χ. υψηλό F1 αλλά χαμηλό Faithfulness λόγω ψευδαισθήσεων), επιβάλλοντας την ισορροπημένη βελτιστοποίηση όλων των συνιστωσών. Σημειώνεται επιπλέον ότι η αξιολόγηση εξαρτάται κρίσιμα από έναν αυτοματοποιημένο ταξινομητή «Δεν Γνωρίζω» (IDK classifier): σωστές αρνήσεις απάντησης σε μη απαντήσιμες στροφές λαμβάνουν τέλειο σκορ (1.0) σε όλους τους άξονες, ενώ εσφαλμένες αρνήσεις ή εσφαλμένες αποδοχές μηδενίζονται — καθιστώντας την πύλη απαντησιμότητας υψηλού ρίσκου συνιστώσα του συστήματος, όπως αναλύεται στην Ενότητα 1.5.4.

%--------------------------------------------------------------
\subsection{Σύνοψη Θεωρητικού Υποβάθρου}
%--------------------------------------------------------------

Συνοψίζοντας, η διασταύρωση των Μεγάλων Γλωσσικών Μοντέλων με τις τεχνικές Νευρωνικής Ανάκτησης Πληροφορίας δημιουργεί ισχυρά υβριδικά συστήματα RAG, η σταθερότητα των οποίων δοκιμάζεται σε πολύστροφα σενάρια διαλόγου. Ο περιορισμός της αστάθειας της ανάκτησης μεταξύ στροφών ---και, σε έναν πραγματικό βοηθό, της σωρευτικής διάδοσης σφαλμάτων--- απαιτεί τον αυστηρό έλεγχο της ποιότητας των ερωτημάτων μέσω πολλαπλών στρατηγικών αναδιατύπωσης, τη χρήση εύρωστων αλγορίθμων σύντηξης κατατάξεων και την εισαγωγή πυλών απαντησιμότητας. Η παρούσα εργασία βασίζεται πάνω σε αυτές τις αρχές για τη διασφάλιση της σταθερότητας της ανάκτησης και της πιστότητας των παραγόμενων αποκρίσεων.
%==============================================================
\section{Ορισμός Προβλήματος}
%==============================================================

Η παρούσα ενότητα εξειδικεύει τον τυπικό μαθηματικό ορισμό του Πολύστροφου RAG και αναλύει τις δομικές προκλήσεις που αυτό συνεπάγεται. Το πρόβλημα προσεγγίζεται ως μια σύνθετη διαδικασία διαδοχικής λήψης αποφάσεων υπό συνθήκες στοχαστικής αβεβαιότητας, η οποία απαιτεί τη σύζευξη δυναμικής ανάκτησης πληροφορίας, παραγωγής φυσικής γλώσσας υπό περιορισμούς (constrained generation) και στοχευμένης εκτίμησης απαντησιμότητας (answerability estimation).

%--------------------------------------------------------------
\subsection{Τυπική Διατύπωση Πολύστροφου RAG}
%--------------------------------------------------------------

Έστω ένας διάλογος εξελισσόμενος σε διαδοχικές στροφές (turns) $i = 1, 2, \dots, t$. Ορίζουμε το διαλογικό ιστορικό $H_t$ μέχρι τη χρονική στιγμή $t$ ως την ακολουθία των προηγούμενων αλληλεπιδράσεων:

\begin{equation}
H_t = \Big( (u_1, \mathcal{D}_1, y_1), (u_2, \mathcal{D}_2, y_2), \dots, (u_{t-1}, \mathcal{D}_{t-1}, y_{t-1}) \Big)
\end{equation}

όπου $u_i$ αντιπροσωπεύει το εκφώνημα (utterance) του χρήστη, $\mathcal{D}_i \subset \mathcal{D}$ το σύνολο των εγγράφων που ανακτήθηκαν και $y_i$ την απόκριση της στροφής $i$. Στο MTRAGEval το ιστορικό αυτό είναι το ιστορικό αναφοράς (reference history) και όχι οι απαντήσεις του ίδιου του συστήματος. Στην τρέχουσα στροφή $t$, ο χρήστης εισάγει ένα νέο, ενδεχομένως μη αυτοτελές ερώτημα $q_t \equiv u_t$.

Με δεδομένο ένα στατικό, καθολικό σώμα γνώσης (knowledge corpus) $\mathcal{D}$, η αρχιτεκτονική Πολύστροφου RAG αποσυντίθεται σε τρεις διασυνδεδεμένες συναρτήσεις, οι οποίες αντιστοιχούν απευθείας στις προδιαγραφές του MTRAGEval:

\paragraph{1. Συνάρτηση Ανάκτησης Συμφραζομένων (Task A):}
Ο μηχανισμός ανάκτησης $\mathcal{R}$ αναλαμβάνει τη χαρτογράφηση του τρέχοντος ερωτήματος και του ιστορικού σε ένα βέλτιστο, διατεταγμένο υποσύνολο $k$ σχετικών αποσπασμάτων (κατά τις προδιαγραφές του διαγωνισμού, $k=10$):
\begin{equation}
\mathcal{R}(q_t, H_t, \mathcal{D}) \rightarrow \mathcal{D}_t = \{d_1, d_2, \dots, d_k\} \subset \mathcal{D}, \quad k=10
\end{equation}

\paragraph{2. Συνάρτηση Τεκμηριωμένης Παραγωγής (Task B):}
Με δεδομένο ένα χρυσό (ground-truth) σύνολο αναφοράς αποδείξεων $\mathcal{D}_t^*$, ο μηχανισμός παραγωγής $\mathcal{G}$ συνθέτει μια πραγματολογικά αγκυρωμένη απόκριση $y_t$, απομονώνοντας πλήρως την ποιότητα παραγωγής από τυχόν σφάλματα ανάκτησης:
\begin{equation}
\mathcal{G}(q_t, H_t, \mathcal{D}_t^*) \rightarrow y_t
\end{equation}

\paragraph{3. Ενοποιημένη End-to-End Παραγωγή (Task C):}
Στο πλήρες σύστημα, η παραγωγή εξαρτάται άμεσα από τα αποτελέσματα του πρώτου σταδίου, γεγονός που εισάγει στοχαστική εξάρτηση και επιτρέπει τη σώρευση σφαλμάτων ανάκτησης στην τελική απόκριση:
\begin{equation}
P(y_t \mid q_t, H_t) = \sum_{\mathcal{D}_t} P(y_t \mid q_t, H_t, \mathcal{D}_t) \cdot P(\mathcal{D}_t \mid q_t, H_t)
\end{equation}

%--------------------------------------------------------------
\subsection{Το Πρόβλημα της Απαντησιμότητας (Answerability Gate)}
%--------------------------------------------------------------

Σε πραγματικά σενάρια συνομιλίας, η πληροφοριακή ανάγκη του χρήστη ενδέχεται να μην καλύπτεται από το διαθέσιμο corpus $\mathcal{D}$. Για τη μοντελοποίηση αυτής της συνθήκης, εισάγουμε μια λανθάνουσα δυαδική μεταβλητή απαντησιμότητας $z_t \in \{0, 1\}$, η οποία ορίζεται ως:

\begin{equation}
z_t =
\begin{cases}
1, & \text{εάν το } \mathcal{D} \text{ περιέχει επαρκή στοιχεία για την πλήρη τεκμηρίωση του } q_t \text{ υπό το } H_t \\
0, & \text{διαφορετικά (unanswerable query)}
\end{cases}
\end{equation}

Το σύστημα οφείλει να εκτιμήσει το $z_t$ μέσω μιας δυαδικής απόφασης $\hat{z}_t \in \{0,1\}$. Στο προτεινόμενο σύστημα, κάθε κριτής $j$ αποφαίνεται για την απαντησιμότητα με βαθμό βεβαιότητας $c_j$, και η στροφή κρίνεται μη απαντήσιμη όταν κάποιος κριτής δηλώνει ανεπάρκεια στοιχείων με βεβαιότητα τουλάχιστον ίση με ένα κατώφλι βαθμονόμησης $\tau$:
\begin{equation}
\hat{z}_t =
\begin{cases}
0, & \text{εάν } \exists j: \text{ο κριτής } j \text{ δηλώνει ανεπάρκεια με } c_j \geq \tau \\
1, & \text{διαφορετικά}
\end{cases}
\end{equation}

Η τελική απόκριση προκύπτει ως:
\begin{equation}
y_t = 
\begin{cases}
\mathcal{G}(q_t, H_t, \mathcal{D}_t), & \text{εάν } \hat{z}_t = 1 \\
\text{«Δεν διαθέτω επαρκείς πληροφορίες για να απαντήσω»}, & \text{εάν } \hat{z}_t = 0
\end{cases}
\end{equation}

Το σφάλμα στην εκτίμηση του $z_t$ παρουσιάζει έντονη ασυμμετρία: τα ψευδώς θετικά αποτελέσματα ($\hat{z}_t=1$ ενώ $z_t=0$) αναγκάζουν το ΜΓΜ να υποπέσει σε πραγματολογικές ψευδαισθήσεις λόγω της τάσης συμμόρφωσης (compliance bias), ενώ τα ψευδώς αρνητικά μειώνουν δραστικά τη χρηστικότητα του συστήματος. Όπως δείχνουμε στην Ενότητα 1.5.4, η ασυμμετρία αυτή δεν είναι θεωρητική μόνο: στο πραγματικό σύστημα εκδηλώνεται ως ισχυρή μεροληψία αποδοχής (acceptance bias) προς το $\hat{z}_t=1$.

%--------------------------------------------------------------
\subsection{Διατύπωση ως Πρόβλημα Ενιαίας Βελτιστοποίησης}
%--------------------------------------------------------------

Η συνολική αρχιτεκτονική στοχεύει στη μεγιστοποίηση της αναμενόμενης χρησιμότητας της απόκρισης υπό μια πολυδιάστατη συνάρτηση κόστους:

\begin{equation}
\mathcal{O}^* = \arg\max_{\mathcal{R}, \mathcal{G}} \mathbb{E}_{(q_t, H_t) \sim \mathcal{X}} \Big[ \mathcal{U}\big(y_t, y_t^*, \mathcal{D}_t, \mathcal{D}_t^*\big) \Big]
\end{equation}

Η συνάρτηση χρησιμότητας $\mathcal{U}$ αναλύεται σε τέσσερις ανεξάρτητους πυλώνες:
\begin{equation}
\mathcal{U} = w_1 \cdot \text{NDCG}(\mathcal{D}_t, \mathcal{D}_t^*) + w_2 \cdot \text{Faithfulness}(y_t, \mathcal{D}_t) + w_3 \cdot \text{Extractive\_F1}(y_t, y_t^*) + w_4 \cdot \Phi(\hat{z}_t, z_t)
\end{equation}
όπου το $\Phi$ ανταμείβει τις σωστά βαθμονομημένες αποφάσεις της πύλης απαντησιμότητας (συμφωνία του $\hat{z}_t$ με το $z_t$)· αφού η $\mathcal{U}$ μεγιστοποιείται, όλοι οι όροι εκφράζονται ως ανταμοιβές. Τα βάρη $w_1, \dots, w_4$ δεν εκτιμώνται αριθμητικά: η διατύπωση είναι εννοιολογική. Λόγω της μη διαφορίσιμης φύσης των μετρικών κειμένου και της ανάκτησης, η βελτιστοποίηση αυτή δεν επιδιώκεται μέσω end-to-end gradient descent, αλλά προσεγγίζεται μέσω σύνθετων agentic pipelines και στοχευμένης μηχανικής προτροπών (prompt engineering), όπως περιγράφεται στην Ενότητα 1.4.

%--------------------------------------------------------------
\subsection{Ιδιαιτερότητες του Πολύστροφου Περιβάλλοντος}
%--------------------------------------------------------------

Το πολύστροφο RAG εισάγει τρεις θεμελιώδεις ακαδημαϊκές προκλήσεις:
\begin{itemize}
    \item \textbf{Χάσμα Λεξιλογίου (Vocabulary Gap):} Τα ερωτήματα μεταγενέστερων στροφών είναι σύντομα και ελλειπτικά. Η απουσία ρητών λέξεων-κλειδιών καθιστά τους κλασικούς ανακτητές (π.χ. BM25) ανίκανους να γεφυρώσουν τη σημασιολογική απόσταση με το corpus.
    \item \textbf{Σωρευτική Διάδοση Σφαλμάτων (Error Cascade):} Σε έναν πραγματικό βοηθό, εάν η συνάρτηση $\mathcal{R}$ αποτύχει στη στροφή $t-1$, ο θόρυβος ενσωματώνεται στο $H_t$ μέσω της απάντησης του συστήματος, και η επόμενη αναδιατύπωση βασίζεται σε μολυσμένα συμφραζόμενα. Το MTRAGEval παρέχει το ιστορικό αναφοράς σε κάθε στροφή, οπότε ο μηχανισμός αυτός δεν ενεργοποιείται στην αξιολόγησή μας· παραμένει όμως η δυσκολία επίλυσης αναφορών σε προηγούμενες στροφές.
    \item \textbf{Μετατόπιση Πλαισίου (Context Drift):} Η δυναμική εναλλαγή θεμάτων εντός του ίδιου διαλόγου απαιτεί από το σύστημα να απομονώνει το σχετικό ιστορικό και να αγνοεί τις ξεπερασμένες πληροφορίες.
\end{itemize}

%--------------------------------------------------------------
\subsection{Κεντρική Υπόθεση Εργασίας}
%--------------------------------------------------------------

Η παρούσα διπλωματική εργασία θεμελιώνεται πάνω στην εξής κεντρική ερευνητική υπόθεση:
\begin{quote}
\textit{«Η συνολική απόδοση ενός end-to-end συστήματος Πολύστροφου RAG φράσσεται άνω (upper-bounded) από τη σταθερότητα της διαδικασίας ανάκτησης (retrieval stability) και τη βαθμονόμηση της πύλης απαντησιμότητας, και όχι από την εγγενή εκφραστική ή συνθετική ικανότητα των σύγχρονων γλωσσικών μοντέλων.»}
\end{quote}
Συνεπώς, η έρευνα εστιάζει στη σταθεροποίηση του $P(\mathcal{D}_t \mid q_t, H_t)$ μέσω ελεγχόμενης ποικιλότητας ερωτημάτων, παρά στην αναζήτηση ισχυρότερων ή μεγαλύτερων γλωσσικών μοντέλων.

%--------------------------------------------------------------
\subsection{Συσχέτιση με το Benchmark MTRAGEval}
%--------------------------------------------------------------

Η αξιολόγηση των προτεινόμενων λύσεων πραγματοποιείται στο απαιτητικό περιβάλλον του SemEval-2026 Task 8, το οποίο αντανακλά τις ανωτέρω προκλήσεις σε ρεαλιστική κλίμακα: (α) η συντριπτική πλειονότητα των στροφών ($87\%$) είναι μη αυτοτελείς, με αναφορές αντωνυμίας να εμφανίζονται στο $97\%$ των μη-πρώτων στροφών· (β) το test set περιλαμβάνει σχεδόν τριπλάσιο ποσοστό μη απαντήσιμων στροφών σε σχέση με το development set· και (γ) παρατηρείται έντονη κατανομική μετατόπιση (distribution shift) μεταξύ development και test set τόσο ως προς την αναλογία unanswerable ερωτημάτων όσο και ως προς την κατανομή των θεματικών πεδίων (domains). Ο Πίνακας~\ref{tab:dev-test-greek} συνοψίζει τα κρίσιμα αυτά μεγέθη.

\begin{table}[h]
\centering
\small
\begin{tabular}{lcc}
\toprule
\textbf{Χαρακτηριστικό} & \textbf{Development} & \textbf{Test} \\
\midrule
Συνομιλίες & 110 & 507 \\
Σύνολο αξιολογούμενων στροφών & 842 & 507 \\
Μέσος αριθμός στροφών/συνομιλία & 7.7 & 1.0 \\
Μη αυτοτελείς στροφές & $87\%$ & — \\
Απαντήσιμες ή μερικώς απαντήσιμες & $92.3\%$ & $74.4\%$ \\
\textbf{Μη απαντήσιμες} & $6.5\%$ & $\mathbf{19.1\%}$ \\
Συνομιλιακές / λοιπές & $1.2\%$ & $6.5\%$ \\
\bottomrule
\end{tabular}
\caption{Σύγκριση development και test set στο MTRAGEval. Η σχεδόν τριπλάσια αύξηση του ποσοστού μη απαντήσιμων στροφών στο test set αποτελεί την κύρια πηγή της κατανομικής μετατόπισης που αναλύεται στην Ενότητα 1.5.5. Στο test set, $377$ στροφές διαθέτουν κρίσεις συνάφειας για το Task~A και $97$ είναι μη απαντήσιμες· η κατηγορία των υπόλοιπων $33$ στροφών ($6.5\%$) δεν μπόρεσε να ανακτηθεί από τα διαθέσιμα αρχεία.}
\label{tab:dev-test-greek}
\end{table}

%--------------------------------------------------------------
\subsection{Ερευνητικά Ερωτήματα}
%--------------------------------------------------------------

Με βάση τη διατύπωση του προβλήματος, η εργασία επιχειρεί να απαντήσει στα εξής τέσσερα ερευνητικά ερωτήματα (Research Questions):
\begin{enumerate}
    \item[\textbf{RQ1:}] Μπορεί η εισαγωγή στοχευμένης ποικιλότητας ερωτημάτων (query diversity) σε έναν ενιαίο ανακτητή να υποκαταστήσει ή να υπερβεί τη σύνθεση ετερογενών μοντέλων ανάκτησης (retriever diversity);
    \item[\textbf{RQ2:}] Σε ποιο βαθμό η αποσύνθεση της παραγωγής με εισαγωγή σταδίου εξαγωγής αποσπασμάτων (verbatim span extraction) συμβάλλει στον περιορισμό των πραγματολογικών ψευδαισθήσεων;
    \item[\textbf{RQ3:}] Ποια είναι η επίδραση της κατανομικής μετατόπισης των unanswerable ερωτημάτων στην ευαισθησία των πυλών απαντησιμότητας;
    \item[\textbf{RQ4:}] Μπορεί η επαναβαθμολόγηση μέσω ΜΓΜ (LLM-based reranking) να προσφέρει μετρήσιμο όφελος έναντι εξειδικευμένων cross-encoder μοντέλων ανάκτησης, όταν η ανάκληση του πρώτου σταδίου ανάκτησης πλησιάζει τον κορεσμό;
\end{enumerate}

%==============================================================
\section{Προτεινόμενο Σύστημα}
%==============================================================

\subsection{Επισκόπηση Αρχιτεκτονικής}
%--------------------------------------------------------------

Η αρχιτεκτονική του προτεινόμενου συστήματος \textbf{AILS-NTUA} σχεδιάστηκε με γνώμονα την αντιμετώπιση των εξαρτήσεων μεταξύ στροφών και της μετατόπισης πλαισίου (context drift) σε πολύστροφα περιβάλλοντα RAG. Η σχεδιαστική φιλοσοφία του συστήματος αποκλίνει από τις παραδοσιακές προσεγγίσεις ετερογενούς σύνθεσης μοντέλων και βασίζεται σε δύο θεμελιώδεις αρχιτεκτονικές αρχές:
\begin{itemize}
    \item \textbf{Ποικιλότητα Ερωτημάτων έναντι Ποικιλότητας Ανακτητών (Query-Diversity-over-Retriever-Diversity):} Αντί για τη χρήση ενός υπολογιστικά δαπανηρού ensemble διαφορετικών ανακτητών (dense και sparse), το σύστημα σταθεροποιεί τον χώρο ανάκτησης εκδίδοντας πολλαπλές, συμπληρωματικές τροχιές συλλογιστικής (prompts) προς έναν ενιαίο, βέλτιστα ευθυγραμμισμένο με το σώμα κειμένου (corpus-aligned) ανακτητή.
    \item \textbf{Καθυστερημένη Ανάληψη Δέσμευσης (Deferred Commitment):} Ο αγωγός παραγωγής αποσυνδέεται από την άμεση σύνθεση απάντησης. Η τελική απόκριση διαμορφώνεται ως μια πολυσταδιακή agentic διαδικασία, η οποία περιλαμβάνει την απομόνωση αποδείξεων σε επίπεδο πρότασης, τη δημιουργία εναλλακτικών υποψηφίων και την τελική επιλογή μέσω ενός συστήματος πολλαπλών εξειδικευμένων κριτών (multi-judge framework).
\end{itemize}

Η ολοκληρωμένη αρχιτεκτονική οργανώνεται σε τρεις λειτουργικές ενότητες που αντιστοιχούν στις προδιαγραφές του διαγωνισμού: τον Αγωγό Ανάκτησης Πολλαπλών Στρατηγικών (\textbf{Task A}, Σχήμα~\ref{fig:taskA-greek}), τον Πολυσταδιακό Αγωγό Παραγωγής (\textbf{Task B}, Σχήμα~\ref{fig:taskB-greek}), και το Ενοποιημένο Σύστημα RAG με Πύλη Απαντησιμότητας (\textbf{Task C}, Σχήμα~\ref{fig:taskC-greek}).

%--------------------------------------------------------------
\subsection{Αγωγός Ανάκτησης Πολλαπλών Στρατηγικών (Task A)}
%--------------------------------------------------------------

Ο αγωγός ανάκτησης εκμεταλλεύεται την ικανότητα των ΜΓΜ να αναλύουν το ιστορικό $H_t$ και να αναδιατυπώνουν το τρέχον ελλιπές ερώτημα $q_t$ υπό πέντε διακριτές, συμπληρωματικές οπτικές γωνίες, με στόχο τον περιορισμό του vocabulary gap:
\begin{enumerate}
    \item \textbf{Ελάχιστη Αναδιατύπωση (Minimal Reformulation):} Εστιάζει αποκλειστικά στην επίλυση αναφορών (coreference resolution) και τη συντακτική συμπλήρωση των ελλείψεων, διατηρώντας τη δομή του αρχικού ερωτήματος.
    \item \textbf{Corpus-Ειδική Αναδιατύπωση (Corpus-Specific):} Εμπλουτίζει το ερώτημα με εξειδικευμένη τεχνική ορολογία και συνώνυμα που ταιριάζουν με το συγκεκριμένο θεματικό πεδίο (π.χ. FiQA, Cloud).
    \item \textbf{Υποθετικό Έγγραφο (HyDE):} Παράγει μια ιδεατή, zero-shot απάντηση στο ερώτημα, μετατοπίζοντας την αναζήτηση από τον χώρο των ερωτήσεων στον χώρο των απαντήσεων του corpus.
    \item \textbf{Αλυσίδα Σκέψης (Chain-of-Thought - CoT):} Παράγει μια σύντομη λογική ανάλυση των επιμέρους πληροφοριακών αναγκών που πρέπει να ικανοποιηθούν για να απαντηθεί το $q_t$.
    \item \textbf{Στόχευση Λέξεων-Κλειδιών (Anchor Keywords):} Εξάγει ονοματοποιημένες οντότητες και στατικά λεκτικά triggers που λειτουργούν ως άγκυρες (anchors) για τον εντοπισμό των εγγράφων.
\end{enumerate}

\begin{figure}[h]
\centering
\begin{tikzpicture}[
  node distance=6mm and 8mm,
  box/.style={draw, rounded corners, align=center, minimum height=8mm, font=\footnotesize, fill=blue!6},
  strat/.style={draw, rounded corners, align=center, minimum height=7mm, font=\scriptsize, fill=orange!8},
  arr/.style={-Stealth, thick}
]
\node[box] (in) {Ερώτημα $q_t$ + Ιστορικό $H_t$};
\node[box, below=of in] (rw) {1. Αναδιατύπωση Πολλαπλών Στρατηγικών};
\node[strat, below left=4mm and -8mm of rw] (s1) {Minimal};
\node[strat, right=2mm of s1] (s2) {Corpus-Sp.};
\node[strat, right=2mm of s2] (s3) {HyDE};
\node[strat, right=2mm of s3] (s4) {CoT};
\node[strat, right=2mm of s4] (s5) {Anchor-KW};
\node[box, below=14mm of rw] (ret) {2. Αραιή Ανάκτηση (ELSER v1) — top-100 ανά ερώτημα};
\node[box, below=of ret] (rer) {3. Υβριδική Επαναβαθμολόγηση: Cohere Rerank v4 + w-RRF};
\node[box, below=of rer] (nrrf) {4. Ιεραρχική Σύντηξη (Nested RRF) — βάρη ανά corpus};
\node[box, below=of nrrf, fill=green!10] (out) {Top-$k$ Αποσπάσματα $\mathcal{D}_t$};
\draw[arr] (in) -- (rw);
\draw[arr] (rw) -- (s1); \draw[arr] (rw) -- (s2); \draw[arr] (rw) -- (s3);
\draw[arr] (rw) -- (s4); \draw[arr] (rw) -- (s5);
\draw[arr] ([yshift=-3mm]s1.south) -- (ret.north -| s1);
\draw[arr] ([yshift=-3mm]s3.south) -- (ret.north -| s3);
\draw[arr] (ret) -- (rer);
\draw[arr] (rer) -- (nrrf);
\draw[arr] (nrrf) -- (out);
\end{tikzpicture}
\caption{Αγωγός ανάκτησης Task A: πέντε συμπληρωματικές αναδιατυπώσεις προωθούνται ανεξάρτητα στον ανακτητή ELSER v1, ακολουθεί υβριδική επαναβαθμολόγηση και τελικά ιεραρχική σύντηξη (Nested RRF).}
\label{fig:taskA-greek}
\end{figure}

\paragraph{Διαδικασία Ανάκτησης και Επαναβαθμολόγησης:}
Κάθε μία από τις πέντε αναδιατυπώσεις προωθείται ανεξάρτητα στον μαθητό αραιό ανακτητή \textbf{ELSER v1}, ο οποίος επιστρέφει τα κορυφαία $N=100$ αποσπάσματα ανά αναδιατύπωση. Στη συνέχεια, τα κορυφαία $60$ έγγραφα κάθε λίστας επεξεργάζονται από ένα cross-encoder μοντέλο επαναβαθμολόγησης (\textbf{Cohere Rerank v4}), του οποίου η βαθμολογία συνδυάζεται με την αρχική κατάταξη του ανακτητή μέσω σταθμισμένης RRF:
\begin{equation}
\text{Score}(d) = \frac{1}{k + r_{E}(d)} + \alpha \cdot \frac{1}{k + r_{R}(d)}
\end{equation}
όπου $r_E$, $r_R$ οι κατατάξεις ανακτητή και reranker αντίστοιχα, $k$ ελέγχει τη φθίνουσα βαρύτητα ανά θέση και $\alpha$ σταθμίζει τα δύο σήματα ($k=60$, $\alpha=0.5$ στη βαθμονομημένη ρύθμιση του συστήματος).

\paragraph{Ιεραρχική Σύντηξη Κατάταξης (Nested RRF):}
Για τη σταθερή ενοποίηση των πέντε επαναβαθμολογημένων κατατάξεων, εισάγεται ένας αλγόριθμος Ιεραρχικής Σύντηξης δύο επιπέδων. Ο σχεδιασμός βασίζεται στην εκτίμηση ότι οι στρατηγικές HyDE, CoT και Anchor-Keyword είναι πιο ασταθείς ανά στροφή και θεματικό πεδίο, ενώ οι Minimal και Corpus-Specific παράγουν πιο σταθερές κατατάξεις· η διακύμανση αυτή δεν μετρήθηκε άμεσα. Στο \textbf{Επίπεδο 1}, οι τρεις διερευνητικές στρατηγικές συγχωνεύονται μέσω RRF για τη διαμόρφωση μιας ενδιάμεσης κατάταξης «Ασθενούς Συναίνεσης» (Weak Consensus). Στο \textbf{Επίπεδο 2}, η κατάταξη αυτή συντίθεται μέσω βεβαρυμμένης RRF με τις δύο σταθερές στρατηγικές, χρησιμοποιώντας βάρη βαθμονομημένα ανά θεματικό πεδίο, παράγοντας την τελική κατάταξη εγγράφων $\mathcal{D}_t$:
\begin{equation}
\text{Score}_{\text{final}}(d) = \sum_{s} w_s^{(c)} \cdot \frac{1}{k^{(c)} + r_s(d)}
\end{equation}
όπου $s$ διατρέχει τις δύο σταθερές στρατηγικές και την ενδιάμεση κατάταξη Weak Consensus, και $w_s^{(c)}$, $k^{(c)}$ είναι παράμετροι εξαρτώμενες από το εκάστοτε corpus $c$.

%--------------------------------------------------------------
\subsection{Πολυσταδιακός Αγωγός Παραγωγής (Task B)}
%--------------------------------------------------------------

Όταν το σύστημα τροφοδοτείται με ένα προκαθορισμένο, έγκυρο σύνολο εγγράφων αναφοράς, η παραγωγή απάντησης εκτελείται μέσω ενός αυστηρά δομημένου agentic pipeline πέντε σταδίων (Σχήμα~\ref{fig:taskB-greek}):

\begin{enumerate}
    \item \textbf{Προκαταρκτικός Έλεγχος Απαντησιμότητας (Answerability Check):} Ένα γρήγορο φίλτρο (gate) εξετάζει αν τα έγγραφα αναφοράς επαρκούν για την κάλυψη του ερωτήματος, αποτρέποντας την ενεργοποίηση των επόμενων, κοστοβόρων σταδίων αν η στροφή είναι unanswerable.
    \item \textbf{Εξαγωγή Αποδεικτικών Εκτάσεων (Evidence Span Extraction):} Το ΜΓΜ αναλύει τα έγγραφα και απομονώνει αυτούσιες (\textit{verbatim}) τις συγκεκριμένες προτάσεις ή φράσεις που περιέχουν την αποδεικτική πληροφορία (έως 8 αποσπάσματα). Αυτό το στάδιο «φιλτράρει» τον περιβάλλοντα θόρυβο και εξαναγκάζει το μοντέλο σε αυστηρή πραγματολογική αγκύρωση (grounding).
    \item \textbf{Παραγωγή Διπλών Υποψηφίων (Dual-Candidate Drafting):} Παράγονται παράλληλα δύο διαφορετικά προσχέδια απαντήσεων από το μοντέλο: ένα \textit{ντετερμινιστικό προσχέδιο} ($T{=}0.0$) με έμφαση στην ακριβή μεταφορά των δεδομένων, και ένα \textit{στοχαστικό προσχέδιο} ($T{=}0.1$) με στόχο τη βελτίωση της συνομιλιακής ροής και της φυσικότητας του λόγου.
    \item \textbf{Επιλογή Απάντησης βάσει Πολλαπλών Κριτών (Multi-Judge Selection):} Τα δύο προσχέδια αξιολογούνται από ένα σχήμα δύο εξειδικευμένων LLM κριτών:
    \begin{itemize}
        \item \textit{Τεχνικός Κριτής (Technical Judge):} Ελέγχει την πραγματολογική ορθότητα και την απουσία αντιφάσεων έναντι των εξαχθέντων spans.
        \item \textit{Κριτής Ικανοποίησης (User Satisfaction Judge):} Καλείται σε στοχαστικό υποσύνολο στροφών ($60\%$) και αξιολογεί αν η απάντηση καλύπτει πλήρως την πληροφοριακή ανάγκη του χρήστη με φυσικό τρόπο.
    \end{itemize}
    Η τελική επιλογή ενσωματώνει επιπλέον έναν όρο \textit{Διαμόρφωσης Εκχύλισης} (Extractiveness Shaping), ο οποίος τιμωρεί τόσο απαντήσεις επιρρεπείς σε ψευδαισθήσεις (πολύ χαμηλό verbatim overlap με τα spans) όσο και υπερβολικά κυριολεκτικές αντιγραφές.
    \item \textbf{Μικροδιορθώσεις και Μορφοποίηση:} Η επιλεγείσα απάντηση υφίσταται τελικό συντακτικό έλεγχο για την αφαίρεση περιττών εκφράσεων εισαγωγής (π.χ. «Βάσει του κειμένου...») και την ευθυγράμμιση του μήκους της με τον στόχο που ορίζεται από τον τύπο ερώτησης.
\end{enumerate}

\begin{figure}[h]
\centering
\begin{tikzpicture}[
  node distance=6mm,
  box/.style={draw, rounded corners, align=center, minimum height=8mm, font=\footnotesize, fill=blue!6},
  cand/.style={draw, rounded corners, align=center, minimum height=7mm, font=\scriptsize, fill=orange!8},
  arr/.style={-Stealth, thick}
]
\node[box] (in) {Έγγραφα Αναφοράς $\mathcal{D}_t^*$};
\node[box, below=of in] (s1) {1. Έλεγχος Απαντησιμότητας};
\node[box, below=of s1] (s2) {2. Εξαγωγή Αποδεικτικών Εκτάσεων (έως 8)};
\node[box, below=of s2] (s3) {3. Παραγωγή Δύο Υποψηφίων};
\node[cand, below left=4mm and -2mm of s3] (ca) {Υποψ. Α ($T{=}0.0$)};
\node[cand, right=4mm of ca] (cb) {Υποψ. Β ($T{=}0.1$)};
\node[box, below=14mm of s3] (s4) {4. Επιλογή: Τεχν. Κριτής + Κριτής Ικανοποίησης + Extractiveness Shaping};
\node[box, below=of s4, fill=green!10] (s5) {5. Μικροδιορθώσεις $\rightarrow$ Τελική Απόκριση $y_t$};
\draw[arr] (in) -- (s1);
\draw[arr] (s1) -- (s2) node[midway, right, font=\tiny] {αν answerable};
\draw[arr] (s2) -- (s3);
\draw[arr] (s3) -- (ca); \draw[arr] (s3) -- (cb);
\draw[arr] (ca) -- (s4); \draw[arr] (cb) -- (s4);
\draw[arr] (s4) -- (s5);
\end{tikzpicture}
\caption{Αγωγός παραγωγής Task B: αλυσίδα πέντε σταδίων που αναβάλλει τη δέσμευση στην τελική απάντηση μέχρι τη συμφωνία πολλαπλών κριτών.}
\label{fig:taskB-greek}
\end{figure}

%--------------------------------------------------------------
\subsection{Ολοκληρωμένο Σύστημα RAG (Task C)}
%--------------------------------------------------------------

Στο end-to-end σενάριο (Task C), οι δύο ανωτέρω αγωγοί ενοποιούνται (Σχήμα~\ref{fig:taskC-greek}). Σε αντίθεση με το Task A, όπου ανακτώνται $k=10$ αποσπάσματα, εδώ το σύστημα περιορίζεται σκόπιμα στα κορυφαία $k=3$ αποσπάσματα — επιλογή που μειώνει τον θόρυβο στην παραγωγή σε συνθήκες αβέβαιης ανάκτησης (βλ. Ενότητα 1.5.4). Για τον περιορισμό των ψευδαισθήσεων, το σύστημα εισάγει μια σύνθετη \textbf{Πύλη Απαντησιμότητας Πολλαπλών Επιπέδων (Multi-Layer Answerability Gate)}.

Αντί της χρήσης ενός μονού κριτή, η απαντησιμότητα αξιολογείται σε τρία διαφορετικά επίπεδα αφαίρεσης:
\begin{itemize}
    \item \textbf{Κριτής Εγγράφου (Document Judge):} Αξιολογεί τη συνολική συνάφεια των ανακτηθέντων κειμένων με το διαλογικό ιστορικό.
    \item \textbf{Κριτής Αποσπάσματος (Span Judge):} Αποφαίνεται για το αν τα εξαχθέντα verbatim spans επαρκούν δομικά για να στηρίξουν μια απάντηση.
    \item \textbf{Κριτής Απάντησης (Answer Judge):} Ελέγχει το τελικό προσχέδιο σε σχέση με το ερώτημα για να εντοπίσει partial ή ελλιπείς πληροφορίες.
\end{itemize}

Οι αποφάσεις των τριών κριτών συνδυάζονται μέσω ενός μηχανισμού \textit{στάθμισης εμπιστοσύνης με δικαίωμα αρνησικυρίας απόκλισης} (confidence-weighted voting with dissenter override): εάν έστω και ένας κριτής δηλώσει με βεβαιότητα τουλάχιστον $\tau{=}0.7$ ότι τα στοιχεία είναι ανεπαρκή, το προσχέδιο απορρίπτεται και επιστρέφεται μια τυποποιημένη άρνηση απάντησης· διαφορετικά, οι αποφάσεις συνδυάζονται με ψηφοφορία σταθμισμένη ως προς τη βεβαιότητα. Επειδή ο Κριτής Απάντησης χρειάζεται προσχέδιο, η εξαγωγή spans και η παραγωγή των δύο υποψηφίων προηγούνται της πύλης· οι απαντήσιμες στροφές συνεχίζουν στα στάδια επιλογής και μικροδιόρθωσης του Task B. Οι τρεις κριτές χρησιμοποιούν το ίδιο πρότυπο προτροπής, αλλά λαμβάνουν διαφορετική είσοδο (αποσπάσματα, spans ή προσχέδιο).

\begin{figure}[ht]
\centering
\begin{tikzpicture}[
  node distance=6mm,
  box/.style={draw, rounded corners, align=center, minimum height=8mm, font=\footnotesize, fill=blue!6},
  jbox/.style={draw, rounded corners, align=center, minimum height=7mm, font=\scriptsize, fill=red!6},
  arr/.style={-Stealth, thick}
]
\node[box] (in) {Ερώτημα $q_t$ + Ιστορικό $H_t$};
\node[box, below=of in] (a) {Αγωγός Task A — top-3 αποσπάσματα};
\node[box, below=of a] (gen) {Εξαγωγή Spans + Παραγωγή 2 Υποψηφίων};
\node[box, below=of gen] (gate) {Πύλη Απαντησιμότητας Πολλαπλών Επιπέδων};
\node[jbox, below left=4mm and -10mm of gate] (j1) {Κριτής Εγγράφου};
\node[jbox, right=2mm of j1] (j2) {Κριτής Αποσπάσματος};
\node[jbox, right=2mm of j2] (j3) {Κριτής Απάντησης};
\node[box, below=16mm of gate] (vote) {Στάθμιση Εμπιστοσύνης + Dissenter Override};
\node[box, below=of vote, fill=green!10] (yes) {Απαντήσιμο $\rightarrow$ Επιλογή + Μικροδιόρθωση};
\node[box, right=8mm of yes, fill=red!10] (no) {Μη απαντήσιμο $\rightarrow$ Τυποποιημένη Άρνηση};
\draw[arr] (in) -- (a); \draw[arr] (a) -- (gen); \draw[arr] (gen) -- (gate);
\draw[arr] (gate) -- (j1); \draw[arr] (gate) -- (j2); \draw[arr] (gate) -- (j3);
\draw[arr] ([yshift=-3mm]j1.south) -- (vote.north -| j1);
\draw[arr] ([yshift=-3mm]j3.south) -- (vote.north -| j3);
\draw[arr] (vote) -- (yes); \draw[arr] (vote) -- (no);
\end{tikzpicture}
\caption{Ενοποιημένος αγωγός Task C: η εξαγωγή spans και η σύνταξη προσχεδίων εκτελούνται στα ανακτημένα αποσπάσματα· στη συνέχεια μια πύλη τριών κριτών είτε επιτρέπει στη στροφή να ολοκληρώσει τα υπόλοιπα στάδια του Task B είτε αντικαθιστά το προσχέδιο με τυποποιημένη άρνηση.}
\label{fig:taskC-greek}
\end{figure}

%--------------------------------------------------------------
\subsection{Σχεδιαστικές Επιλογές και Συνεισφορά}
%--------------------------------------------------------------

Η αρχιτεκτονική συνεισφορά του συστήματος AILS-NTUA συνοψίζεται στις ακόλουθες καινοτόμες σχεδιαστικές επιλογές:
\begin{itemize}
    \item \textbf{Αντικατάσταση των Ετερογενών Ensembles:} Ένδειξη, στο development set, ότι η στοχευμένη υπολογιστική επένδυση στην ποικιλότητα των ερωτημάτων (query diversity) υπερτερεί της χρήσης πολλών διαφορετικών μοντέλων ανάκτησης.
    \item \textbf{Αλγόριθμος Nested RRF:} Εισαγωγή ιεραρχικής δομής στη σύντηξη κατατάξεων, η οποία συγκεντρώνει τις διερευνητικές στρατηγικές σε μία ενδιάμεση κατάταξη πριν τη συνδυάσει με τις συντηρητικές.
    \item \textbf{Αποσύνθεση Παραγωγής μέσω Spans:} Η μετάβαση από το text-to-text RAG στο span-guided generation, γεγονός που ελαχιστοποιεί τον θόρυβο εισόδου στον generator.
    \item \textbf{Βαθμονομημένο Multi-Judge Voting:} Η ανάπτυξη μιας πύλης απαντησιμότητας πολλών επιπέδων, σχεδιασμένης να περιορίζει τις ψευδαισθήσεις πριν από την τελική δέσμευση της απάντησης (inference commitment).
\end{itemize}
%==============================================================
\section{Πειραματικά Αποτελέσματα}
%==============================================================

Η παρούσα ενότητα παρουσιάζει την ποσοτική και ποιοτική αξιολόγηση του συστήματος AILS-NTUA. Αναλύεται η πειραματική διάταξη, παρατίθενται οι επίσημες επιδόσεις στα leaderboards του SemEval-2026 Task 8, και πραγματοποιείται εκτενής ανάλυση συνιστωσών (ablation studies) για τον προσδιορισμό της συνεισφοράς κάθε αρχιτεκτονικής επιλογής στα Tasks A, B και C.

%--------------------------------------------------------------
\subsection{Πειραματική Διάταξη}
%--------------------------------------------------------------

Η αξιολόγηση του συστήματος πραγματοποιήθηκε στο επίσημο benchmark του MTRAGEval (SemEval-2026 Task 8). Το σύνολο δεδομένων περιλαμβάνει $110$ πολύστροφες συνομιλίες που εκτείνονται σε συνολικά $842$ στροφές (turns) και καλύπτουν τέσσερα ετερογενή θεματικά πεδία (domains): \textit{FiQA} (οικονομικά), \textit{ClapNQ} (γενικές γνώσεις), \textit{Govt} (κυβερνητικά έγγραφα) και \textit{Cloud} (τεχνική υποστήριξη). Η ετερογένεια αυτή είναι σκόπιμη και δομικά σημαντική: όπως φαίνεται στον Πίνακα~\ref{tab:corpus-greek}, τα τέσσερα corpora διαφέρουν σημαντικά ως προς τη μέση πυκνότητα αποσπασμάτων ανά έγγραφο — από ατομικά έγγραφα-αναρτήσεις στο FiQA έως εκτενή εγκυκλοπαιδικά λήμματα στο ClapNQ — στρεσάροντας διαφορετικές πτυχές του μηχανισμού ανάκτησης. Η αξιολόγηση του Task A (Passage Retrieval) περιορίζεται στις $777$ στροφές που χαρακτηρίζονται ως απαντήσιμες (answerable) ή μερικώς απαντήσιμες (partially answerable).

\begin{table}[ht]
\centering
\small
\begin{tabular}{lcccc}
\toprule
\textbf{Χαρακτηριστικό} & \textbf{ClapNQ} & \textbf{FiQA} & \textbf{Govt} & \textbf{Cloud} \\
\midrule
Τύπος πεδίου & Εγκυκλοπαιδικό & Οικονομικό φόρουμ & Κυβερνητικά κείμενα & Τεχνική τεκμηρίωση \\
Σύνολο εγγράφων & 4{,}293 & 57{,}638 & 7{,}661 & 8{,}578 \\
Σύνολο αποσπασμάτων & 183{,}408 & 61{,}022 & 49{,}607 & 72{,}442 \\
Μέσος όρος αποσπ./έγγρ. & 42.7 & 1.1 & 6.5 & 8.4 \\
Κύρια πηγή & Wikipedia & StackExchange & Κυβ. ιστότοποι & IBM Cloud Docs \\
\bottomrule
\end{tabular}
\caption{Στατιστικά δομικής ετερογένειας του corpus MTRAG ανά θεματικό πεδίο.}
\label{tab:corpus-greek}
\end{table}

\paragraph{Υπερπαράμετροι και Ρύθμιση Συστήματος:}
Η βαθμονόμηση όλων των υπερπαραμέτρων πραγματοποιήθηκε αποκλειστικά στο development set. Στον αγωγό ανάκτησης, κάθε μία από τις $5$ LLM-based αναδιατυπώσεις εκδίδει ερώτημα προς τον ανακτητή \textbf{ELSER v1}, επιστρέφοντας τα κορυφαία $100$ αποσπάσματα κειμένου. Η επαναβαθμολόγηση εκτελείται μέσω του μοντέλου \textbf{Cohere Rerank v4} για τα $60$ επικρατέστερα έγγραφα, με σταθερά RRF $k=60$, ενώ η ιεραρχική συγχώνευση (Nested RRF) χρησιμοποιεί σταθερές εξομάλυνσης $k_{\text{internal}}=40$ στο Επίπεδο~1 και $k_{\text{final}}$ από $20$ έως $60$ ανά θεματικό πεδίο στο Επίπεδο~2.

\paragraph{Υπολογιστική Δρομολόγηση Μοντέλων (Model Routing):}
Για τη διαχείριση του οικονομικού κόστους των API, υιοθετήθηκε μια αρχιτεκτονική δρομολόγησης βάσει της κρισιμότητας του ρόλου (role-based model routing):
\begin{itemize}
    \item \textbf{DeepSeek-V3.2:} Χρησιμοποιείται για τις εργασίες υψηλής δομικής ακρίβειας, όπως η εξαγωγή των verbatim αποδεικτικών εκτάσεων (span extraction) και η υλοποίηση του Τεχνικού Κριτή (Technical Judge).
    \item \textbf{GPT-4o:} Αξιοποιείται για την Παραγωγή Υποψηφίων Απαντήσεων (Dual-Candidate Drafting) και για τους τρεις κριτές της πύλης απαντησιμότητας του Task C.
    \item \textbf{GPT-4o-mini:} Χρησιμοποιείται για υπολογιστικά ελαφριές εργασίες, όπως ο προκαταρκτικός έλεγχος απαντησιμότητας, ο Κριτής Ικανοποίησης Χρήστη και οι τελικές μικροδιορθώσεις.
\end{itemize}

%--------------------------------------------------------------
\subsection{Αποτελέσματα Ανάκτησης — Task A}
%--------------------------------------------------------------

\paragraph{Συγκριτική Ανάλυση Μοντέλων Ανάκτησης:}
Η αξιολόγηση εννέα ανακτητών (λεξιλογικών, πυκνών και learned sparse) στο development set, \textit{χωρίς} αναδιατύπωση ερωτήματος, κατέδειξε την ξεκάθαρη υπεροχή των learned sparse αναπαραστάσεων. Ο ανακτητής ELSER v1 πέτυχε $Recall@5: 0.483$ και $nDCG@5: 0.444$, υπερτερώντας του ισχυρότερου πυκνού ανακτητή (Cohere Embed v4, $R@5{=}0.397$) κατά $+8.6$ ποσοστιαίες μονάδες, και του κλασικού λεξιλογικού BM25 ($R@5{=}0.153$) κατά $+33.0$ ποσοστιαίες μονάδες στα ίδια $777$ ερωτήματα. Μέρος του πλεονεκτήματος αυτού μπορεί να οφείλεται στο ότι οι κρίσεις συνάφειας του MTRAG παρήχθησαν κυρίως από αποσπάσματα του ELSER. Το αποτέλεσμα είναι συνεπές με την υπόθεση ότι η επέκταση όρων (term expansion) σε επίπεδο corpus μειώνει το vocabulary gap χωρίς τον σημασιολογικό θόρυβο των cross-domain dense embeddings, αλλά δεν μπορεί να διαχωριστεί πλήρως από την παραπάνω μεροληψία των κρίσεων.

\paragraph{Ανάλυση Συνεισφοράς Αναδιατυπώσεων (Query Reformulation Ablation):}
Ο Πίνακας~\ref{tab:rewrite-ablation-greek} καταγράφει τη σωρευτική επίδραση της σταδιακής προσθήκης των πέντε στρατηγικών αναδιατύπωσης πάνω στον ανακτητή ELSER v1. Ο πλήρης συνδυασμός και των πέντε στρατηγικών υπό το αλγόριθμο Nested RRF αυξάνει το $Recall@5$ κατά \textbf{$+25.7\%$} σε σχέση με τη βασική γραμμή χωρίς αναδιατύπωση. Σε επίπεδο μεμονωμένης στρατηγικής, η \textit{Corpus-Specific} αναδεικνύεται ως η ισχυρότερη επιλογή στα πεδία ClapNQ ($R@5{=}0.637$) και FiQA ($R@5{=}0.491$), η \textit{Chain-of-Thought} υπερτερεί στο Govt ($R@5{=}0.580$), ενώ η \textit{Minimal} παραμένει η πλέον αποδοτική στο τεχνικό domain Cloud ($R@5{=}0.463$) — το HyDE, αντίθετα, αποδεικνύεται συστηματικά η ασθενέστερη μεμονωμένη στρατηγική σε όλα τα domains, χρήσιμη όμως ως συμπληρωματικό σήμα εντός του σχήματος Nested RRF. Ο πλήρης συνδυασμός ($R@5{=}0.607$) υπερβαίνει κατά $6.6$ μονάδες ακόμη και την καλύτερη μεμονωμένη στρατηγική· επειδή όμως περιλαμβάνει και την επαναβαθμολόγηση, η διαφορά είναι συνεπής με τη \textit{συμπληρωματικότητα} των στρατηγικών χωρίς να την απομονώνει.

\begin{table}[ht]
\centering
\small
\begin{tabular}{lcc}
\toprule
\textbf{Διαμόρφωση Ανάκτησης} & \textbf{Recall@5} & \textbf{Σχετική Βελτίωση ($\Delta$)} \\
\midrule
Μη ενισχυμένη βάση (αρχικό ερώτημα) & 0.483 & — \\
+ Ελάχιστη Αναδιατύπωση (Minimal) & 0.527 & $+9.1\%$ \\
+ Corpus-Ειδική Επέκταση & 0.558 & $+15.5\%$ \\
+ Αλυσίδα Σκέψης (CoT) & 0.573 & $+18.6\%$ \\
+ Υποθετικό Έγγραφο (HyDE) & 0.584 & $+20.9\%$ \\
+ Στόχευση Λέξεων-Κλειδιών & 0.591 & $+22.4\%$ \\
+ Σύντηξη Nested RRF (πλήρες σύστημα) & \textbf{0.607} & \textbf{$+25.7\%$} \\
\bottomrule
\end{tabular}
\caption{Σωρευτική ανάλυση συνεισφοράς (ablation) των πέντε στρατηγικών αναδιατύπωσης ερωτήματος στο development set. Οι γραμμές είναι σωρευτικές, οπότε κάθε αύξηση εξαρτάται από τη σειρά προσθήκης· ο πίνακας δεν απομονώνει τη συνεισφορά του reranker ή της ιεραρχικής σύντηξης.}
\label{tab:rewrite-ablation-greek}
\end{table}

\paragraph{Το Παράδοξο της Σύνθεσης Ανακτητών (Retriever Ensemble Paradox):}
Ένα από τα πιο σημαντικά ευρήματα της εργασίας είναι ότι η σύνθεση ετερογενών μοντέλων (π.χ. ELSER + BM25 + Dense μέσω επίπεδης RRF) υποβαθμίζει την ακρίβεια κορυφής (top-rank precision), παρότι βελτιώνει την οριακή ανάκληση σε μεγάλο βάθος. Όπως φαίνεται στον Πίνακα~\ref{tab:ensemble-paradox-greek}, στο πλήρως βελτιστοποιημένο στάδιο του αγωγού (ELSER με τις πέντε αναδιατυπώσεις), η προσθήκη όλων των εναλλακτικών ανακτητών αυξάνει το $Recall@100$ από $0.901$ σε $0.926$, ταυτόχρονα όμως \textit{μειώνει} το $Recall@5$ από $0.607$ σε $0.569$. Μια εξήγηση είναι ότι τα μοναδικά χρυσά έγγραφα που εισάγουν οι εναλλακτικοί ανακτητές εμφανίζονται σε μέση θέση από $27{,}5$ έως $54{,}2$ μετά τη σύντηξη ($39{,}5$ για όλους μαζί), κάτω δηλαδή από το όριο top-$10$, ενώ ταυτόχρονα η σύντηξη μετατοπίζει οριακά σχετικά αποτελέσματα του ELSER εκτός των κορυφαίων θέσεων. Η προτεινόμενη προσέγγιση \textit{ενιαίου ανακτητή με ποικιλότητα ερωτημάτων} αποφεύγει αυτή την παγίδα, καθώς όλα τα σήματα παραμένουν ευθυγραμμισμένα με τον ίδιο χώρο βαθμολόγησης. Σημειώνεται ότι οι κρίσεις συνάφειας του MTRAG παρήχθησαν κυρίως από αποσπάσματα που ανέκτησε το ELSER· αποσπάσματα που φέρνουν μόνο οι υπόλοιποι ανακτητές είναι πιθανότερο να μην έχουν κριθεί και να μετρώνται ως άσχετα. Η μεροληψία αυτή μπορεί από μόνη της να συμβάλλει στο παράδοξο, και δεν ποσοτικοποιήθηκε.

\begin{table}[ht]
\centering
\small
\begin{tabular}{lccc}
\toprule
\textbf{Διαμόρφωση (τελικό στάδιο, 5 αναδιατυπώσεις)} & \textbf{R@5} & \textbf{R@100} & \textbf{Μέση θέση μοναδικών χρυσών εγγρ.} \\
\midrule
ELSER v1 (μόνο) & \textbf{0.607} & 0.901 & — \\
+ Cohere Embed (dense) & 0.582 $\downarrow$ & 0.925 $\uparrow$ & 37.0 \\
+ SPLADE & 0.592 $\downarrow$ & 0.919 $\uparrow$ & 54.2 \\
+ BM25 & 0.536 $\downarrow$ & 0.915 $\uparrow$ & 27.5 \\
+ Όλοι οι ανακτητές & 0.569 $\downarrow$ & \textbf{0.926} $\uparrow$ & 39.5 \\
\bottomrule
\end{tabular}
\caption{Το παράδοξο σύνθεσης ανακτητών: η προσθήκη εναλλακτικών ανακτητών βελτιώνει το $R@100$ αλλά υποβαθμίζει συστηματικά το $R@5$.}
\label{tab:ensemble-paradox-greek}
\end{table}

\paragraph{Σταθερότητα ανά Στροφή Διαλόγου:}
Η απόδοση ανάκτησης εξαρτάται έντονα από τη θέση της στροφής (Σχήμα~\ref{fig:turn-depth-domain}). Χωρίς αναδιατύπωση, το $Recall@5$ πέφτει από $0.879$ στην πρώτη στροφή σε $0.428$ στις στροφές $6+$ (σχετική πτώση $51\%$)· με το πλήρες σύστημα η πτώση περιορίζεται στο $39\%$ (από $0.890$ σε $0.540$). Η πτώση δεν είναι σταδιακή: συγκεντρώνεται κυρίως μεταξύ της πρώτης στροφής, που είναι εξ ορισμού αυτοτελής, και της δεύτερης, ενώ στη συνέχεια οι καμπύλες μένουν σχεδόν σταθερές. Συνεπώς η επίδραση της θέσης οφείλεται κυρίως στο αν το ερώτημα είναι αυτοτελές ή εξαρτάται από προηγούμενες στροφές (αναφορές αντωνυμίας στο $97\%$ των μη-πρώτων στροφών, συγκριτικές ερωτήσεις-συνέχειες στο $67\%$). Η αναδιατύπωση περιορίζει, χωρίς να εξαλείφει, το χάσμα αυτό.

%--------------------------------------------------------------
\subsection{Αποτελέσματα Τεκμηριωμένης Παραγωγής — Task B}
%--------------------------------------------------------------

\paragraph{Ablation Analysis του Agentic Pipeline:}
Ο Πίνακας~\ref{tab:gen-ablation-greek} καταγράφει τη σωρευτική συνεισφορά κάθε σταδίου του πολυσταδιακού αγωγού παραγωγής στο Task B. Η σταδιακή ενσωμάτωση των συνιστωσών οδηγεί σε συνολική αύξηση του Αρμονικού Μέσου (HM) κατά \textbf{$+8.9$ μονάδες} ($0.654 \rightarrow 0.743$) έναντι μίας απλής, μονοβηματικής (greedy) παραγωγής.

\begin{table}[ht]
\centering
\small
\begin{tabular}{lc}
\toprule
\textbf{Διαμόρφωση Αγωγού} & \textbf{Αρμονικός Μέσος (HM)} \\
\midrule
Βάση: Μονοβηματική Greedy Παραγωγή ($T{=}0.0$) & 0.654 \\
+ Εξαγωγή Αποδεικτικών Εκτάσεων (Verbatim) & 0.688 \\
+ Παραγωγή Δύο Υποψηφίων (τυχαία επιλογή) & 0.706 \\
+ Τεχνικός Κριτής & 0.725 \\
+ Κριτής Ικανοποίησης Χρήστη ($60\%$ δειγματοληψία) & 0.738 \\
+ Τελικές Μικροδιορθώσεις & \textbf{0.743} \\
\bottomrule
\end{tabular}
\caption{Ανάλυση συνεισφοράς (ablation) του αγωγού παραγωγής Task B στο development set.}
\label{tab:gen-ablation-greek}
\end{table}

\paragraph{Η Κρισιμότητα του Evidence Span Extraction:}
Η εισαγωγή του σταδίου εξαγωγής αποδεικτικών εκτάσεων αναδεικνύεται ως η \textit{μεμονωμένα} μεγαλύτερη βελτίωση του αγωγού ($+0.034$ HM, Πίνακας~\ref{tab:gen-ablation-greek}), μαζί με αύξηση κατά $+0.041$ στη μετρική πιστότητας χωρίς αναφορά $RL_F$. Η verbatim απομόνωση των σχετικών προτάσεων περιορίζει την τάση των μοντέλων να ενσωματώνουν άσχετες πληροφορίες (distractors) από το υπόλοιπο σώμα των εγγράφων, εμποδίζοντας έτσι τον generator να «δραπετεύσει» από τα παρεχόμενα στοιχεία.

\paragraph{Αποτίμηση του Σχήματος Πολλαπλών Κριτών:}
Η εισαγωγή των δύο LLM κριτών (Technical, User Satisfaction) σε συνδυασμό με τον όρο Διαμόρφωσης Εκχύλισης (Extractiveness Shaping) στοχεύει στην εξισορρόπηση της πραγματολογικής ακρίβειας με τη συνομιλιακή φυσικότητα, και κάθε στάδιο προσθέτει θετική συνεισφορά στο ablation του Πίνακα~\ref{tab:gen-ablation-greek}. Ο όρος εκχύλισης ποινικοποιεί προσχέδια που είτε απομακρύνονται υπερβολικά από τα εξαχθέντα spans (κίνδυνος ψευδαίσθησης) είτε τα αντιγράφουν σχεδόν αυτολεξεί, ευνοώντας επικάλυψη 4-γραμμάτων εντός της ζώνης $[0.28, 0.38]$, που επιλέχθηκε στο development set (Πίνακας~\ref{tab:extractiveness-band-en}).

\paragraph{Οικονομική Αποδοτικότητα μέσω Routing:}
Η δρομολόγηση μοντέλων (DeepSeek-V3.2 για εξαγωγή/τεχνική κρίση, GPT-4o για τη σύνθεση, GPT-4o-mini για ελαφριές εργασίες) είναι σχεδιαστική επιλογή με κίνητρο το κόστος. Με ονομαστικές τιμές API, ο αγωγός του Task B κοστίζει κατά εκτίμηση $\$0{,}010$--$0{,}012$ ανά στροφή, από τα οποία τα $\$0{,}009$ αντιστοιχούν στις δύο κλήσεις του GPT-4o, έναντι εκτιμώμενων $\$0{,}024$ αν όλα τα στάδια χρησιμοποιούσαν GPT-4o, δηλαδή περίπου $2$--$2{,}4$ φορές χαμηλότερο κόστος.\footnote{Το δημοσιευμένο άρθρο~\citep{athanasiou2026ailsntua} αναφέρει $\$0{,}008$ ανά στροφή και $3\times$ μείωση. Το άθροισμα των εκτιμήσεων ανά στάδιο, με προσμέτρηση και των δύο κλήσεων παραγωγής, δίνει $\$0{,}010$--$0{,}012$, τιμή που υιοθετείται εδώ.} Δεν διατηρήθηκε ελεγχόμενη σύγκριση ποιότητας με ομοιόμορφη χρήση ενός μοντέλου, οπότε δεν διατυπώνεται ισχυρισμός για την επίδραση της δρομολόγησης στην ποιότητα.

%--------------------------------------------------------------
\subsection{Η Απαντησιμότητα ως Καθοριστικός Παράγοντας --- Task C}
%--------------------------------------------------------------

Η end-to-end αξιολόγηση (Task C) παρουσιάζει απότομη πτώση σε σχέση με την παραγωγή με αποσπάσματα αναφοράς ($HM$: $0.7698 \rightarrow 0.5409$ στο test set). Σε συμφωνία με το δημοσιευμένο άρθρο~\citep{athanasiou2026ailsntua}, θεωρούμε τη βαθμονόμηση της απαντησιμότητας καθοριστικό παράγοντα υπό την κατανομή του test set: η πύλη αποδέχεται τις περισσότερες μη απαντήσιμες στροφές ακόμη και στο development set και βαθμονομήθηκε σε ποσοστό μη απαντήσιμων $6.5\%$, ενώ στο test set αντιμετωπίζει $19.1\%$. Η ανάκτηση παραμένει πάντως σημαντικός παράγοντας. Στο Task C ο generator λαμβάνει μόνο τα τρία πρώτα αποσπάσματα, οπότε το αν τα αποδεικτικά στοιχεία φτάνουν σε αυτόν καθορίζεται από την ακρίβεια κορυφής ($Recall@5{=}0.607$ στο development set, Πίνακας~\ref{tab:ensemble-paradox-greek}) και όχι από την ανάκληση σε βάθος ($Recall@100{=}0.901$). Επιπλέον, τα σφάλματα ανάκτησης στις κορυφαίες θέσεις αντιστοιχούν στις περισσότερες αποτυχίες του development set που εξετάστηκαν χειροκίνητα (Πίνακας~\ref{tab:error-decomp-greek}), σε ένα σύνολο όπου οι μη απαντήσιμες στροφές είναι σπάνιες. Η ανάλυση των επίσημων σκορ του test set ανά κατηγορία απαντησιμότητας θα ποσοτικοποιούσε τις δύο συνεισφορές. Όπως φαίνεται στον Πίνακα~\ref{tab:answerability-greek}, καμία διαμόρφωση πύλης δεν ξεπερνά το $25\%$ $F1$ στην ανίχνευση μη απαντήσιμων στροφών στο πλήρες development set. Το προτεινόμενο σχήμα πολλαπλών κριτών πετυχαίνει $95.8\%$ ανάκληση στις απαντήσιμες στροφές αλλά μόλις $21.8\%$ στις μη απαντήσιμες, δηλαδή $12$ από τις $55$ (95\% διάστημα Wilson $12{,}9$--$34{,}4\%$), ένδειξη ισχυρής μεροληψίας αποδοχής (acceptance bias). Η ποιοτική εξέταση των αποφάσεων της πύλης δείχνει ότι οι κριτές Εγγράφου και Αποσπάσματος συχνά επισημαίνουν ανεπαρκή στοιχεία, ενώ ο Κριτής Απάντησης τείνει να αποδέχεται προσχέδια που παραμένουν εύλογα, δίνοντας προτεραιότητα στην ευλογοφάνεια έναντι της επάρκειας των αποδείξεων.

\begin{table}[ht]
\centering
\small
\begin{tabular}{lccc}
\toprule
\textbf{Μέθοδος} & \textbf{Ανάκληση Απαντ.} & \textbf{Ανάκληση Μη Απαντ.} & \textbf{F1 Μη Απαντ.} \\
\midrule
Πάντα Απάντηση (χωρίς πύλη) & $100.0\%$ (777/777) & $0.0\%$ (0/55) & $0.0\%$ \\
Αγωγός + Κανόνας Επαλήθευσης & $92.3\%$ (717/777) & $27.3\%$ (15/55) & $23.1\%$ \\
\textbf{Πολλαπλοί Κριτές (δικό μας)} & $95.8\%$ (744/777) & $21.8\%$ (12/55) & $24.0\%$ \\
\bottomrule
\end{tabular}
\caption{Μετρικές ανίχνευσης απαντησιμότητας στο development set ($832$ στροφές: $777$ απαντήσιμες ή μερικώς απαντήσιμες και $55$ μη απαντήσιμες· εξαιρούνται οι $10$ συνομιλιακές στροφές). Καμία διαμόρφωση που αξιολογήθηκε δεν ξεπερνά το $25\%$ $F1$ στις μη απαντήσιμες στροφές του πλήρους development set.}
\label{tab:answerability-greek}
\end{table}

\paragraph{Ποσοτική Κατανομή Σφαλμάτων και Εξαρτήσεις μεταξύ Στροφών:}
Ο συγγραφέας εξέτασε χειροκίνητα $50$ περιπτώσεις αποτυχίας του development set και απέδωσε σε καθεμία μία κύρια αιτία (Πίνακας~\ref{tab:error-decomp-greek}). Οι περιπτώσεις δεν επιλέχθηκαν με τεκμηριωμένη τυχαία διαδικασία και η κατηγοριοποίηση έγινε από έναν σχολιαστή, οπότε τα ποσοστά είναι ενδεικτικά.

\begin{table}[ht]
\centering
\small
\begin{tabular}{lcc}
\toprule
\textbf{Αιτία Σφάλματος} & \textbf{Ποσοστό} & \textbf{Κύριο Πεδίο} \\
\midrule
Σφάλμα ανάκτησης (περιεχόμενο τρέχουσας στροφής) & $38\%$ & FiQA \\
Σφάλμα ανάκτησης (ανεπίλυτη εξάρτηση από προηγούμενες στροφές) & $34\%$ & Cloud \\
Εσφαλμένη εκτίμηση Απαντησιμότητας & $18\%$ & Όλα \\
Καθαρό σφάλμα παραγωγής (σωστά αποσπάσματα) & $10\%$ & ClapNQ \\
\bottomrule
\end{tabular}
\caption{Κύρια αιτία σε $50$ περιπτώσεις αποτυχίας του Task C που εξετάστηκαν χειροκίνητα (development set, ένας σχολιαστής). Οι αιτίες που σχετίζονται με την ανάκτηση αντιστοιχούν στο $72\%$.}
\label{tab:error-decomp-greek}
\end{table}

Οι αιτίες που σχετίζονται με την ανάκτηση αντιστοιχούν συνολικά στο $72\%$ των περιπτώσεων και η εσφαλμένη εκτίμηση απαντησιμότητας στο $18\%$. Η δεύτερη κατηγορία ανάκτησης ($34\%$) χρειάζεται προσεκτική ανάγνωση: επειδή το MTRAGEval παρέχει το ιστορικό αναφοράς και όχι τις προηγούμενες απαντήσεις του συστήματος, ένα σφάλμα του συστήματος στη στροφή $t-1$ δεν μπορεί να περάσει στο $H_t$. Η κατηγορία $34\%$ αφορά επομένως στροφές των οποίων το ερώτημα εξαρτάται από προηγούμενες στροφές και η αναδιατύπωση δεν επέλυσε την εξάρτηση, με αποτέλεσμα να αποτύχει η ανάκτηση στη στροφή $t$ παρότι το ιστορικό είναι σωστό. Οι περιπτώσεις αυτές υποδεικνύουν την επίλυση συναναφοράς και έλλειψης ως σημαντική εναπομένουσα αδυναμία της ανάκτησης, σε συμφωνία με το χάσμα μεταξύ αυτοτελών και μη αυτοτελών ερωτημάτων.

%--------------------------------------------------------------
\subsection{Επίσημα Αποτελέσματα Leaderboard}
%--------------------------------------------------------------

Στην επίσημη διεθνή αξιολόγηση του SemEval-2026 Task 8, η αρχιτεκτονική \textbf{AILS-NTUA} κατέλαβε την 1η θέση στο Task A και τη 2η στο Task B (Πίνακας~\ref{tab:leaderboard-greek}):
\begin{itemize}
    \item \textbf{1η Θέση στο Task A (Passage Retrieval):} Το σύστημα πέτυχε την κορυφαία επίδοση με \textbf{$nDCG@5: 0.5776$}, σημειώνοντας διαφορά \textbf{$+20.5\%$} έναντι της ισχυρότερης επίσημης γραμμής βάσης (baseline) του διαγωνισμού, μεταξύ $38$ συμμετοχών.
    \item \textbf{2η Θέση στο Task B (Grounded Generation):} Το σύστημα κατέλαβε τη δεύτερη θέση με Αρμονικό Μέσο \textbf{$HM: 0.7698$} ($RB_{alg}{=}0.633$, $RL_F{=}0.897$, $RB_{llm}{=}0.832$), σε συμφωνία με τα ευρήματα για το span extraction και το multi-judge framework, μεταξύ $26$ συμμετοχών.
    \item \textbf{11η Θέση στο Task C (End-to-End RAG):} Το σύστημα πέτυχε $HM: 0.5409$ μεταξύ $29$ συμμετοχών, μόλις οριακά πάνω από την ισχυρότερη γραμμή βάσης των διοργανωτών ($0.5366$, $+0.004$). Η πτώση από το $HM{=}0.7698$ του Task B οφείλεται κυρίως στη μετρική $RB_{alg}$ ($0.633 \rightarrow 0.400$) και συμφωνεί με την ανάλυση της Ενότητας 1.5.4: η πύλη βαθμονομήθηκε σε development set με ποσοστό μη απαντήσιμων στροφών περίπου ίσο με το ένα τρίτο εκείνου του test set ($6.5\%$ έναντι $19.1\%$, Πίνακας~\ref{tab:dev-test-greek}), οπότε το κατώφλι $\tau{=}0.7$ δεν ταιριάζει στις συνθήκες του test set. Συμβάλλουν επίσης τα σφάλματα ανάκτησης στις κορυφαίες θέσεις, αφού στον generator φτάνουν μόνο τρία αποσπάσματα. Το αν ένα χαμηλότερο κατώφλι (π.χ. $\tau \approx 0.6$) θα απέδιδε καλύτερα στο test set αποτελεί υπόθεση που δεν μπορεί να ελεγχθεί με μία μόνο υποβολή· στο development set, το $\tau{=}0.6$ αποδίδει ελαφρώς χειρότερα από το $\tau{=}0.7$.
\end{itemize}

\begin{table}[ht]
\centering
\small
\begin{tabular}{llcc}
\toprule
\textbf{Task} & \textbf{Μετρική} & \textbf{Επίδοση} & \textbf{Κατάταξη} \\
\midrule
A — Ανάκτηση & nDCG@5 & 0.5776 & 1/38 \\
B — Παραγωγή & HM & 0.7698 & 2/26 \\
C — End-to-End RAG & HM & 0.5409 & 11/29 \\
\bottomrule
\end{tabular}
\caption{Επίσημα αποτελέσματα του AILS-NTUA στο test set του SemEval-2026 Task 8.}
\label{tab:leaderboard-greek}
\end{table}
%==============================================================
\section{Συμπεράσματα και Μελλοντική Εργασία}
%==============================================================

Η παρούσα ενότητα ολοκληρώνει την εκτεταμένη περίληψη της διπλωματικής εργασίας. Πραγματοποιείται μια συνθετική αποτίμηση των αρχιτεκτονικών συνεισφορών του συστήματος AILS-NTUA, δίνονται ρητές απαντήσεις στα θεμελιώδη ερευνητικά ερωτήματα που καθοδήγησαν την έρευνα, αναλύονται οι εγγενείς περιορισμοί της παρούσας υλοποίησης και χαράσσονται οι μελλοντικές ερευνητικές κατευθύνσεις για τον χώρο των πολύστροφων συστημάτων RAG.

%--------------------------------------------------------------
\subsection{Σύνοψη Συνεισφορών}
%--------------------------------------------------------------

Η παρούσα εργασία διερεύνησε συστηματικά το πρόβλημα της πολύστροφης Ανάκτησης-Επαυξημένης Παραγωγής (Multi-Turn RAG) υπό το πρίσμα της σταθερότητας ανάκτησης συμφραζομένων (context-aware retrieval stability) και της αξιόπιστης πραγματολογικής αγκύρωσης αποδείξεων. Σε αντίθεση με τις συμβατικές προσεγγίσεις της βιβλιογραφίας που περιορίζονται στη βελτιστοποίηση μεμονωμένων σταδίων, αναπτύχθηκε μια ενοποιημένη αρχιτεκτονική που διαχωρίζει λειτουργικά τη διατήρηση του διαλογικού ιστορικού, την ανακάλυψη αποδείξεων και την παραγωγή απάντησης υπό περιορισμούς.

Οι κύριες ερευνητικές και τεχνολογικές συνεισφορές της εργασίας συνοψίζονται ως εξής:
\begin{itemize}
    \item \textbf{Αναδιατύπωση ως Πρωτεύον Σήμα Ανάκτησης (Query Diversity):} Στο development set και με βάση τις επίσημες κρίσεις συνάφειας του benchmark, η εισαγωγή στοχευμένης ποικιλότητας ερωτημάτων μέσω πέντε συμπληρωματικών LLM προτροπών πάνω σε έναν ενιαίο, corpus-ευθυγραμμισμένο ανακτητή (ELSER v1) υπερτερεί της σύνθεσης ετερογενών μοντέλων ανάκτησης (retriever ensembles) (Ενότητα 1.5.2).
    \item \textbf{Αλγόριθμος Ιεραρχικής Σύντηξης Κατατάξεων (Nested RRF):} Σχεδιάστηκε μια στρατηγική σύντηξης δύο επιπέδων, η οποία συγκεντρώνει τις διερευνητικές στρατηγικές (HyDE, Chain-of-Thought, Anchor-Keyword) σε μια ενδιάμεση κατάταξη πριν τη συνδυάσει με τις συντηρητικές. Η πλήρης διάταξη ανάκτησης περιορίζει την πτώση της απόδοσης μετά την πρώτη στροφή από $51\%$ σε $39\%$, χωρίς όμως να έχει απομονωθεί η συνεισφορά της ιεραρχικής σύντηξης αυτής καθαυτής.
    \item \textbf{Διύλιση Πλαισίου μέσω Evidence Span Extraction:} Η εισαγωγή ενός ενδιάμεσου σταδίου verbatim εξαγωγής αποδεικτικών προτάσεων πριν από την παραγωγή περιορίζει την επίδραση των σημασιολογικών παραπλανητών (distractors) και αποτελεί τη μεμονωμένα μεγαλύτερη συνεισφορά στον αγωγό παραγωγής ($+0.034$ HM στο development set).
    \item \textbf{Agentic Επιλογή Απάντησης Δύο Κριτών:} Ένα σχήμα εξειδικευμένων LLM κριτών (Technical, User Satisfaction) σε συνδυασμό με όρο Διαμόρφωσης Εκχύλισης (Extractiveness Shaping) στοχεύει στην εξισορρόπηση της πραγματολογικής πιστότητας με τη φυσικότητα της συνομιλιακής ροής.
    \item \textbf{Ανάλυση του Task C:} Η πτώση Task B$\rightarrow$C στο test set ($HM: 0.7698 \rightarrow 0.5409$) υποδεικνύει ως καθοριστικό παράγοντα τη βαθμονόμηση της πύλης απαντησιμότητας υπό σχεδόν τριπλάσιο ποσοστό μη απαντήσιμων ερωτημάτων, ενώ τα σφάλματα ανάκτησης στις κορυφαίες θέσεις συμβάλλουν σημαντικά· στο development set, όπου οι μη απαντήσιμες στροφές είναι σπάνιες, τα σφάλματα ανάκτησης στις κορυφαίες θέσεις εξηγούν το $72\%$ των αποτυχιών που εξετάστηκαν χειροκίνητα· καμία διαμόρφωση πύλης δεν ξεπερνά το $25\%$ F1 στο πλήρες development set (Πίνακας~\ref{tab:answerability-greek}).
\end{itemize}

%--------------------------------------------------------------
\subsection{Απαντήσεις στα Ερευνητικά Ερωτήματα}
%--------------------------------------------------------------

Τα τέσσερα κεντρικά ερευνητικά ερωτήματα (Research Questions) που διατυπώθηκαν στην Ενότητα 1.3 απαντώνται βάσει των πειραματικών δεδομένων ως εξής:
\begin{itemize}
    \item[\textbf{RQ1:}] \textit{Μπορεί η εισαγωγή στοχευμένης ποικιλότητας ερωτημάτων να υποκαταστήσει τη σύνθεση ετερογενών ανακτητών;} \\
    \textbf{Απάντηση:} Ναι, με βάση τις επίσημες κρίσεις συνάφειας. Η ποικιλότητα ερωτημάτων σε έναν ενιαίο ανακτητή υπερτερεί στο development set και αντιστοιχεί στην 1η θέση στο Task A ($nDCG@5: 0.5776$). Αντίθετα, η προσθήκη εναλλακτικών ανακτητών μειώνει το $Recall@5$, καθώς τα χρυσά έγγραφα που φέρνουν εμφανίζονται σε μέση θέση από $27{,}5$ έως $54{,}2$ μετά τη σύντηξη (Πίνακας~\ref{tab:ensemble-paradox-greek}), κάτω δηλαδή από το όριο top-$10$. Μέρος του φαινομένου μπορεί να οφείλεται στη μεροληψία των κρίσεων συνάφειας υπέρ του ELSER.
    \item[\textbf{RQ2:}] \textit{Σε ποιο βαθμό η εξαγωγή verbatim αποδεικτικών εκτάσεων συμβάλλει στον περιορισμό των ψευδαισθήσεων;} \\
    \textbf{Απάντηση:} Στο development set, η εξαγωγή spans αποτελεί τη μεγαλύτερη μεμονωμένη βελτίωση στον αγωγό παραγωγής ($+0.034$ Αρμονικό Μέσο), μαζί με αύξηση $+0.041$ στη μετρική πιστότητας χωρίς αναφορά ($RL_F$). Επειδή μια μεταβολή σε μία μόνο συνιστώσα αυξάνει τον αρμονικό μέσο τριών μετρικών μόνο κατά ένα κλάσμα της, το κέρδος $+0.034$ στον HM δείχνει ότι βελτιώθηκαν και οι μετρικές με αναφορά· η ευφράδεια δεν μετρήθηκε χωριστά. Ο περιορισμός της εισόδου του generator σε επαληθεύσιμες προτάσεις περιορίζει τον χώρο στον οποίο μπορούν να εισαχθούν αστήρικτοι ισχυρισμοί· το υψηλό $RL_F$ του test set ($0.897$) είναι συνεπές με αυτό το εύρημα.
    \item[\textbf{RQ3:}] \textit{Ποια είναι η επίδραση της κατανομικής μετατόπισης των unanswerable ερωτημάτων στην ευαισθησία των πυλών;} \\
    \textbf{Απάντηση:} Ουσιαστική. Η βαθμονόμηση της πύλης στο development set (κατώφλι $\tau{=}0.7$, ποσοστό unanswerable $6.5\%$) είναι πιθανότατα ακατάλληλη για το test set, όπου το ποσοστό σχεδόν τριπλασιάζεται ($19.1\%$), και η πτώση Task B$\rightarrow$C προέρχεται κυρίως από τη μετρική $RB_{alg}$ ($0.633 \rightarrow 0.400$). Τα στοιχεία αυτά αναδεικνύουν τη βαθμονόμηση της απαντησιμότητας ως καθοριστικό παράγοντα της end-to-end απόδοσης υπό την κατανομή του test set, αν και συμβάλλουν και τα σφάλματα ανάκτησης στις κορυφαίες θέσεις· οι δύο επιδράσεις δεν μπορούν να διαχωριστούν χωρίς ανάλυση των σκορ του test set ανά κατηγορία. Το σχήμα πολλαπλών κριτών βελτιώνει τον μονό κριτή στο development set, αλλά αφήνει ισχυρό acceptance bias.
    \item[\textbf{RQ4:}] \textit{Μπορεί η επαναβαθμολόγηση μέσω ΜΓΜ (LLM-based reranking) να προσφέρει όφελος έναντι εξειδικευμένων cross-encoder μοντέλων όταν η ανάκληση του πρώτου σταδίου είναι κοντά στον κορεσμό;} \\
    \textbf{Απάντηση:} Όχι ουσιαστικά. Καμία από τις τρεις μορφές LLM reranking (pointwise, listwise, generation-based) δεν βελτίωσε τη σταθμισμένη σύντηξη ανακτητή και cross-encoder όταν το $Recall@100$ ξεπερνά το $0.90$· ως βοηθητικό σήμα με μικρό βάρος δίνουν το πολύ οριακά κέρδη (έως $+2$ μονάδες $Recall@5$ ανά θεματικό πεδίο), που δεν δικαιολογούν το πρόσθετο κόστος.
\end{itemize}

%--------------------------------------------------------------
\subsection{Περιορισμοί του Συστήματος}
%--------------------------------------------------------------

Παρά τις υψηλές επιδόσεις στα Tasks A και B, το σύστημα και η αξιολόγησή του έχουν συγκεκριμένους περιορισμούς:
\begin{enumerate}
    \item \textbf{Ρύθμιση και αναφορά στο ίδιο σύνολο:} Όλες οι υπερπαράμετροι ρυθμίστηκαν στο development set, στο οποίο αναφέρονται και όλα τα ablations, χωρίς ξεχωριστό τμήμα ελέγχου ή cross-validation ανά συνομιλία· τα αποτελέσματα του development set είναι επομένως πιθανώς αισιόδοξα. Κάθε διαμόρφωση εκτελέστηκε μία φορά, οπότε μικρές διαφορές (της τάξης του $0.01$ HM ή λιγότερο) δεν πρέπει να υπερερμηνεύονται.
    \item \textbf{Στατικότητα Κατωφλίων Απόφασης:} Η Πύλη Απαντησιμότητας βασίζεται σε σταθερό κατώφλι εμπιστοσύνης ($\tau=0.7$), βαθμονομημένο αποκλειστικά στο development set. Το κατώφλι αυτό, μαζί με τη ζώνη εκχύλισης ($[0.28, 0.38]$) και τα βάρη του Nested RRF ανά corpus, δεν επανα-βαθμονομήθηκαν για το test set. Η αναντιστοιχία αυτή είναι πιθανότατα σημαντικός παράγοντας του χάσματος Task B$\rightarrow$C, μαζί με τα σφάλματα ανάκτησης στις κορυφαίες θέσεις· το αν ένα άλλο κατώφλι θα απέδιδε καλύτερα στο test set δεν μπορεί να επαληθευτεί με μία μόνο υποβολή.
    \item \textbf{Ιστορικό αναφοράς στην αξιολόγηση:} Το MTRAGEval παρέχει σε κάθε στροφή το ιστορικό αναφοράς, οπότε η αξιολόγηση δεν ελέγχει την ανθεκτικότητα του συστήματος στα δικά του προηγούμενα σφάλματα, όπως θα συνέβαινε σε έναν πραγματικό βοηθό.
    \item \textbf{Μεροληψία των κρίσεων συνάφειας:} Οι κρίσεις συνάφειας του MTRAG παρήχθησαν κυρίως από αποσπάσματα του ELSER, γεγονός που μπορεί να ευνοεί το ELSER έναντι άλλων ανακτητών και να συμβάλλει στο παράδοξο σύνθεσης ανακτητών.
    \item \textbf{Εξάρτηση από Κλειστά API (Proprietary Models):} Η αρχιτεκτονική βασίζεται στο GPT-4o και στο DeepSeek-V3.2, γεγονός που αυξάνει το λειτουργικό κόστος και μειώνει την αναπαραγωγιμότητα, λόγω της έλλειψης πρόσβασης στα βάρη των μοντέλων. Κανένα στάδιο του αγωγού δεν έχει υποστεί fine-tuning στα δεδομένα MTRAG· όλα τα κέρδη προέρχονται από μηχανική προτροπών (prompt engineering) και αρχιτεκτονικές επιλογές.
    \item \textbf{Υπολογιστική Πολυπλοκότητα Pipeline:} Η πολυσταδιακή φύση του αγωγού (5 αναδιατυπώσεις, cross-encoder reranking, πύλη απαντησιμότητας τριών κριτών) αυξάνει τη δομική πολυπλοκότητα και τον χρόνο απόκρισης, ο οποίος δεν μετρήθηκε, και απαιτεί πρόσθετη μηχανική ενορχήστρωσης (orchestration) για χρήση σε κλίμακα.
\end{enumerate}

%--------------------------------------------------------------
\subsection{Μελλοντικές Ερευνητικές Κατευθύνσεις}
%--------------------------------------------------------------

Η παρούσα διπλωματική εργασία ανοίγει τις ακόλουθες ερευνητικές κατευθύνσεις:
\begin{itemize}
    \item \textbf{Προσαρμοστική Βαθμονόμηση (Adaptive Thresholding):} Δυναμικοί μηχανισμοί απαντησιμότητας, όπου το κατώφλι $\tau$ μεταβάλλεται βάσει της εκτίμησης της εντροπίας των ανακτηθέντων εγγράφων και της ανίχνευσης μετατοπίσεων κατανομής, αντί να παραμένει στατικό. Επειδή η βαθμονόμηση της απαντησιμότητας είναι ο καθοριστικός παράγοντας υπό την κατανομή του test set, πρόκειται για κατεύθυνση υψηλής αναμενόμενης αξίας· η ανάλυση των επίσημων σκορ του test set ανά κατηγορία απαντησιμότητας θα βοηθούσε να ποσοτικοποιηθεί το αναμενόμενο κέρδος.
    \item \textbf{Ισχυρότερες γραμμές βάσης για την πύλη:} Ταξινομητές επάρκειας βασισμένοι σε NLI, ταξινομητές πάνω στα σκορ ανάκτησης και κατώφλια βαθμονομημένα ανά θεματικό πεδίο, που δεν αξιολογήθηκαν στην παρούσα εργασία.
    \item \textbf{Ανθεκτική επίλυση εξαρτήσεων μεταξύ στροφών:} Η δεύτερη μεγαλύτερη κατηγορία σφαλμάτων είναι η ανεπίλυτη εξάρτηση από προηγούμενες στροφές ($34\%$). Καλύτερη επίλυση αναφορών και ελλείψεων ---π.χ. με ρητή παρακολούθηση οντοτήτων--- στοχεύει άμεσα σε αυτήν. Σε πραγματικά συστήματα, όπου το ιστορικό περιέχει τις απαντήσεις του ίδιου του συστήματος, μηχανισμοί ελέγχου και αναίρεσης προηγούμενων βημάτων (backtracking) θα μπορούσαν επιπλέον να περιορίσουν το error cascade.
    \item \textbf{Fine-Tuning Ανοιχτών Μοντέλων Συλλογιστικής (Open-Weight Reasoning):} Εκπαίδευση εξειδικευμένων open-weight μοντέλων (π.χ. Llama-3-Reasoning, DeepSeek-R1) μέσω απόσταξης γνώσης (knowledge distillation) από το προτεινόμενο pipeline, με στόχο τη μείωση του κόστους και την άρση της εξάρτησης από κλειστά API.
    \item \textbf{Ενισχυτική Μάθηση (RL) για Query Formulation:} Εφαρμογή αλγορίθμων ενισχυτικής μάθησης (π.χ. PPO) για την εκπαίδευση του query rewriter, με συνάρτηση ανταμοιβής (reward function) απευθείας το nDCG@5 του ανακτητή.
\end{itemize}

%--------------------------------------------------------------
\subsection{Τελικό Συμπέρασμα}
%--------------------------------------------------------------

Συμπερασματικά, η παρούσα εργασία υποδεικνύει ότι η επιτυχής υλοποίηση συστημάτων Πολύστροφου RAG δεν εξαρτάται αποκλειστικά από τη χρήση ισχυρότερων γλωσσικών μοντέλων. Η μετατόπιση της έμφασης από την απλή κλιμάκωση (scaling) προς τον προσεκτικό αρχιτεκτονικό σχεδιασμό, τη \textit{σταθερότητα της ανάκτησης} και την \textit{αυστηρή βαθμονόμηση των αποφάσεων} αποτελεί μια υποσχόμενη κατεύθυνση για συστήματα υψηλής αξιοπιστίας και διαφάνειας. Τα ευρήματα της εργασίας στηρίζονται στις αναλύσεις του development set και στα επίσημα αποτελέσματα, με την 1η και τη 2η θέση στα Tasks A και B· η υστέρηση στο Task C αναδεικνύει τη \textit{βαθμονόμηση της απαντησιμότητας}, δηλαδή την απόφαση απάντησης ή αποχής υπό μετατόπιση κατανομής, ως το καθοριστικό ανοιχτό πρόβλημα, με την ακρίβεια κορυφής της ανάκτησης να παραμένει σημαντικός δευτερεύων περιορισμός.
